\documentclass[11pt]{article}

\usepackage[a4paper,margin=22mm,headheight=14pt]{geometry}
\usepackage{fontspec}
\usepackage{xeCJK}
\setmainfont{DejaVu Serif}
\setsansfont{DejaVu Sans}
\setmonofont{DejaVu Sans Mono}
\setCJKmainfont{Noto Serif CJK KR}
\setCJKsansfont{Noto Sans CJK KR}
\setCJKmonofont{Noto Sans Mono CJK KR}

\usepackage{amsmath,amssymb,bm,mathtools}
\usepackage{booktabs,longtable,tabularx,array}
\usepackage{enumitem}
\usepackage{xcolor}
\usepackage{hyperref}
\usepackage{fancyhdr}
\usepackage{microtype}
\usepackage{listings}
\usepackage[most]{tcolorbox}

\definecolor{CoreBlue}{HTML}{1F4E79}
\definecolor{CoreLight}{HTML}{EAF3F8}
\definecolor{DeferOrange}{HTML}{B45F06}
\definecolor{DeferLight}{HTML}{FFF2CC}
\definecolor{RejectRed}{HTML}{9C0006}
\definecolor{RejectLight}{HTML}{FCE8E6}
\definecolor{SafeGreen}{HTML}{38761D}
\definecolor{SafeLight}{HTML}{EAF4E2}

\hypersetup{
  colorlinks=true,
  linkcolor=CoreBlue,
  citecolor=CoreBlue,
  urlcolor=CoreBlue,
  pdftitle={Topology-Distilled Selective Global-Local Wind Dynamics for Static 3D Gaussian Thin Surfaces},
  pdfauthor={Wind3DGS Project}
}

\pagestyle{fancy}
\fancyhf{}
\lhead{\small Topology-Distilled Selective Global--Local Wind Dynamics}
\rhead{\small Minimal V0}
\cfoot{\thepage}

\setlength{\parindent}{0pt}
\setlength{\parskip}{0.45em}
\setlist[itemize]{leftmargin=1.5em,itemsep=0.15em,topsep=0.25em}
\setlist[enumerate]{leftmargin=1.8em,itemsep=0.2em,topsep=0.25em}
\renewcommand{\arraystretch}{1.18}
\newcolumntype{L}[1]{>{\raggedright\arraybackslash}p{#1}}
\newcommand{\code}[1]{\texttt{\small #1}}
\newcommand{\core}[1]{\textcolor{CoreBlue}{\textbf{#1}}}
\newcommand{\defer}[1]{\textcolor{DeferOrange}{\textbf{#1}}}
\newcommand{\reject}[1]{\textcolor{RejectRed}{\textbf{#1}}}
\newcommand{\safe}[1]{\textcolor{SafeGreen}{\textbf{#1}}}

\lstset{
  basicstyle=\ttfamily\small,
  breaklines=true,
  frame=single,
  backgroundcolor=\color{black!2},
  columns=fullflexible
}

\title{\textbf{Topology-Distilled Selective Global--Local Wind Dynamics}\\
for Static 3D Gaussian Thin Surfaces\\[0.35em]
\large Minimal V0 Canonical Research Sketch}
\author{Wind3DGS Project}
\date{2026-08-15}

\begin{document}
\maketitle

\begin{tcolorbox}[colback=CoreLight,colframe=CoreBlue,title=Canonical method/contract freeze,fonttitle=\bfseries]
이 문서는 첫 CG 논문의 method equation과 implementation contract를 소유하고 동결하는
Minimal V0 canonical source다. 논문의 최종 claim과 연구 경로는 이 동결과 구분하며,
Gate A/B/C evidence가 모이는 V0-05에서만 별도로 동결한다.
구현 checklist는 이 문서의 stable label을 참조하는 파생 실행 문서이며,
독립적인 두 번째 수식 source가 아니다.
목표는 engineering-grade thin-shell simulation이 아니라
\emph{visually plausible, temporally stable, controllable wind animation}과
실측 quality--latency 이득이다.
이전 full design의 objective normal-curvature closure, patch-force KKT/Schur,
exact reaction ledger, same-frame corrector와 confidence calibration은
design archive로 동결하며 V0 구현 요구사항이 아니다.
\end{tcolorbox}

\tableofcontents
\newpage

\section{연구 질문과 검증 대상 기여 계층}

\subsection{연구 질문}

독립적으로 재구성된 static 3D Gaussian Splatting(3DGS) asset
\cite{kerbl2023gaussian}에서 explicit target triangle mesh를 복원하지 않고,
사용자가 지정한 재료와 부착 조건 아래 바람에 반응하는 얇은 물체를
빠르고 직접 렌더 가능한 형태로 애니메이션할 수 있는가?

핵심 가설은 전체 reduced rank를 무조건 늘리는 대신, Global ROM이 놓치는
국소 응답의 complementary Local modes만 선택하면 strong Global-rank 및
same-anchor solver보다 새로운 quality--latency Pareto 영역을 만들 수 있다는 것이다.
이는 learned gate의 존재가 아니라 선택적 complementary execution 자체의 가설이다.

\subsection{기여 계층}

\begin{enumerate}
  \item \textbf{MC1 --- topology-distilled sparse physical scaffold.}
  Independent static GS에서 oriented sparse anchors, local relation, area/mass ownership과
  Global/patch-Local basis를 만들되 target에서 explicit triangle-mesh reconstruction을 요구하지 않는다.
  Material과 attachment는 appearance로 추론하지 않고 setup metadata로 받는다.
  \item \textbf{MC2 --- selective complementary Global--Local dynamics.}
  Global basis의 rank를 전체적으로 늘리지 않고 현재 load에 필요한 patch complement만 선택하여
  하나의 coupled reduced system으로 적분한다.
  \item \textbf{SYS --- directly renderable Gaussian transport and scaling evidence.}
  Anchor dynamics를 Gaussian mean/covariance로 한 번 transport하고,
  simulation anchor 수와 rendering Gaussian 수를 분리해 측정한다.
\end{enumerate}

Structural backend, quasi-steady aero, selector score와 affine transport는 enabling component다.
각 component 자체를 독립 novelty로 주장하지 않는다.
기존 Gaussian physics/animation 연구 \cite{xie2024physgaussian,yan2026gaussianswaying,
lei2026diffwind,ma2026fastphysgs}와의 차이는 위 MC1/MC2 및 실측 Pareto에서 입증한다.
방어할 교차점은 independent static GS의 versioned sparse scaffold, static-GS-derived Global space에
\(M\)-orthogonal한 patch complement의 wind-load-conditioned fixed-cardinality execution,
그리고 anchor-to-Gaussian direct transport를 하나의 measured system으로 닫는 결합이다.
Mesh-free ROM, local enrichment, Gaussian aerodynamics 또는 bending triplet 각각의 최초성을 주장하지 않는다.

\subsection{명시적 비주장}

\begin{itemize}
  \item Engineering-grade stress, force, torque, support reaction 또는 material identification
  \item Full nonlinear shell accuracy, nonzero Poisson ratio, anisotropy와 pre-curved shell coupling
  \item Wake, vortex shedding, aerodynamic history 또는 two-way FSI
  \item Self-contact, tearing, plastic crease, detached/free flight
  \item Exact energy conservation, exact modal recovery 또는 exact error-controlled adaptivity
  \item Fully mesh-free physics 또는 arbitrary thin-object generalization
  \item End-to-end p95 측정 전의 real-time/interactive claim
\end{itemize}

\section{대상 범위와 I/O contract}

\subsection{정량 범위}

첫 구현과 Gate 판정의 정량 core는 flat-rest, single-layer, attached thin surface에 한정한다.
Core geometry는 cantilever strip과 rectangular flag이며, uniform isotropic material preset
2--3개와 hard edge attachment 하나를 먼저 사용한다. Flat-rest acceptance는 각 same-layer
support의 식~\eqref{eq:minimal-flat-rest-planarity} normalized rest-normal residual이 development에서 동결한
\(\tau_{\mathrm{planar}}\) 이하일 것을 요구한다. Irregular flat sampling에서는 아래
opposite-triplet descriptor \(d_{ia}^0\)가 정확히 0이 아닐 수 있으므로 이를 0으로 덮어쓰지 않고,
기존 tangential interpolation residual과 planarity cap을 함께 적용한다.
Steady, abrupt start/stop, sinusoidal gust와 traveling gust만 core wind program으로 둔다.

Gate A/B/C 통과 후에만 shallow-rest/pre-curved surface, triangular pennant, curtain strip,
compliant attachment와 real captured flag를 제한적인 qualitative transfer로 추가한다.
이때도 frozen V0 backend를 그대로 사용하며 새 shell solver를 추가하지 않고 정량 core claim에
포함하지 않는다. Paper-like sheet와 leaf/petal도 qualitative transfer이며,
plastic crease, anisotropy와 stem mechanics를 주장하지 않는다.

\subsection{Setup 입력}

\begin{equation}
  \mathcal I_{\mathrm{setup}}
  =
  \left(
    \mathcal G^0,\;
    s_{\mathrm{metric}}\ \text{or}\ \vartheta_{\mathrm{nd}},\;
    \vartheta_{\mathrm{mat}},\;
    \vartheta_{\mathrm{aero}},\;
    \mathcal C_{\mathrm{attach}}
  \right).
  \label{eq:minimal-setup-input}
\end{equation}

\(\mathcal G^0\)는 static Gaussian means, covariances와 appearance다.
Unit branch는 \(\mathsf u\in\{\mathrm{SI},\mathrm{nd}\}\) 중 정확히 하나다.
SI branch는 \(s_{\mathrm{metric}}>0\)와 직접 지정한 surface mass density
\(\rho_A\,[\mathrm{kg}/\mathrm{m}^2]\)를 사용한다. Volumetric density와 thickness를 이미 아는 경우에만
\(\rho_A=\rho_s h\)로 환산하며, \(h\)를 appearance에서 추론하지 않는다.
Native GS length unit에서 SI package로 들어가는 metric conversion은
\begin{equation}
  x^{\mathrm{SI}}=s_{\mathrm{metric}}x^{\mathrm{native}},\qquad
  \Sigma^{\mathrm{SI}}=s_{\mathrm{metric}}^2\Sigma^{\mathrm{native}},\qquad
  a^{\mathrm{SI}}=s_{\mathrm{metric}}^2a^{\mathrm{native}}
  \label{eq:minimal-si-metric-conversion}
\end{equation}
으로 고정한다. 즉 Gaussian covariance와 anchor area를 length와 같은 비율로 변환하지 않는다.
실제 SI scale이 없으면
\(\vartheta_{\mathrm{nd}}=(L_0,M_0,T_0)\), \(L_0,M_0,T_0>0\)와
별도의 dimensionless material/aero preset을 사용한다. SI payload를 명시적인
\((L_0,M_0,T_0)\)로 nondimensional internal representation에 옮길 때와 SI/nd parity fixture에는
다음 typed map만 사용한다.
\begin{equation}
\begin{aligned}
  \widehat x&=x/L_0,&
  \widehat t&=t/T_0,&
  \widehat{\Delta t}&=\Delta t/T_0,\\
  \widehat m&=m/M_0,&
  \widehat a&=a/L_0^2,&
  \widehat\Sigma&=\Sigma/L_0^2,\\
  \widehat\rho_A&=\rho_A L_0^2/M_0,&
  \widehat\kappa_{\mathrm{str}}&=\kappa_{\mathrm{str}}T_0^2/M_0,&
  \widehat\kappa_{\mathrm{bend}}&=\kappa_{\mathrm{bend}}T_0^2/(M_0L_0^2),\\
  \widehat\rho_{\mathrm{air}}&=\rho_{\mathrm{air}}L_0^3/M_0,&
  \widehat W&=WT_0/L_0,&
  \widehat{\partial_tW}&=(\partial_tW)T_0^2/L_0,\\
  \widehat{\nabla_xW}&=T_0\nabla_xW,&
  \widehat f&=fT_0^2/(M_0L_0),&
  \widehat f_{\max}&=f_{\max}T_0^2/(M_0L_0).
\end{aligned}
  \label{eq:minimal-si-nd-typed-map}
\end{equation}
여기서 \(\rho_A,\kappa_{\mathrm{str}},\kappa_{\mathrm{bend}},\rho_{\mathrm{air}}\)의 SI 차원은 각각
\(M/L^2,M/T^2,ML^2/T^2,M/L^3\)이고,
\(W,\partial_tW,\nabla_xW,f\)의 차원은 각각 \(L/T,L/T^2,1/T,ML/T^2\)다.
Nondimensional material payload는
\(\widehat\vartheta_{\mathrm{mat}}
=(\widehat\rho_A,\widehat\kappa_{\mathrm{str}},
\widehat\kappa_{\mathrm{bend}},\widehat\alpha,\widehat\beta)\)를 직접 저장한다.
절대 mass가 없으면 \(\widehat M_{\mathrm{asset}}=1\)과
\(\widehat\rho_A=1/\sum_i\widehat a_i^0\)를 기본 convention으로 사용한다.
SI와 nondimensional preset에는 서로 다른 type tag와 package hash를 부여하고,
한 package 또는 run에서 섞이면 fail-fast한다.
\(\vartheta_{\mathrm{mat}}
=(\rho_A,\kappa_{\mathrm{str}},\kappa_{\mathrm{bend}},\alpha,\beta)\)는
unit-typed surface mass density, stretching stiffness, bending rigidity와 Rayleigh damping preset이며,
SI Rayleigh law \(C=\alpha M+\beta K\)에서
\([\alpha]=\mathrm{s}^{-1}\), \([\beta]=\mathrm{s}\)다.
SI preset을 \((L_0,M_0,T_0)\) convention으로 nondimensionalize하면
\begin{equation}
  \widehat\alpha=\alpha T_0,\qquad
  \widehat\beta=\frac{\beta}{T_0}
  \label{eq:minimal-rayleigh-unit-map}
\end{equation}
을 저장한다. 처음부터 nondimensional preset을 지정한 경우에는 이 두 값 자체가 typed input이며,
SI coefficient와 같은 field에 암묵적으로 섞지 않는다.
Length, area, covariance, velocity, mass, force와 energy bound/epsilon은 각각
\(L_0,L_0^2,L_0^2,L_0/T_0,M_0,M_0L_0/T_0^2,M_0L_0^2/T_0^2\)로 나누고,
dimensionless rank/condition/rate threshold는 변환하지 않는다. 선택한 internal unit mode,
\((L_0,M_0,T_0)\)의 값과 소유자, 변환된 clamp/eigenvalue/epsilon 전체를 package hash에 넣는다.
SI internal mode를 선택한 package는 존재하지 않는 scale triple을 암묵적으로 만들어 위 map을 적용하지 않는다.
\(\vartheta_{\mathrm{aero}}=(\rho_{\mathrm{air}},C_n,C_\tau,\text{clamp})\)는
run 전 동결하는 quasi-steady aerodynamic preset이다.
\(\mathcal C_{\mathrm{attach}}\)는 setup에서 고정하는 hard attachment layout이다.
Runtime 중 material 또는 attachment topology를 바꾸지 않는다.
두 branch는 package build에서 하나의 선택된 internal unit representation으로 정규화한 뒤
동일한 operator/solver code path를 사용한다.

\subsection{Frame 입력과 출력}

Frame 입력은 prescribed one-way wind field \(W(x,t)\), timestep \(\Delta t\)와 이전 state다.
출력은 deformed Gaussian means/covariances, standard 3DGS renderer 입력과
stage별 synchronized timing trace다.

Canonical run의 operating envelope는
\begin{equation}
  \mathcal E_{\mathrm{op}}
  =
  \left(
    \Delta t,\;N_{\mathrm{sub}}=1,\;W_{\max},\;\dot W_{\max},\;G_{W,\max},\;
    f_{\max},\;\mathcal B_{\mathrm{area/affine/cov}}
  \right)
  \label{eq:minimal-operating-envelope}
\end{equation}
로 package/run manifest에 저장한다. Core wind는
\(\|W\|\leq W_{\max}\), \(\|\partial_tW\|\leq\dot W_{\max}\),
\(\|\nabla_xW\|\leq G_{W,\max}\) 안에서만 정량 주장하며,
force와 area/affine/covariance clamp bound 및 clamp-rate cap도 함께 동결한다.
\(f_{\max}>0\)는 선택한 internal unit branch에 맞게 typed된 positive scalar이며,
anchor별 \(\|f_{\mathrm{aero},i}\|_2\)의 hard bound다.
Canonical package attempt는 결과를 보기 전에 등록한 independent-GS reconstruction과 setup metadata의
조합이고, canonical sequence attempt는 그 package에 대해 미리 등록한 material/wind/seed/horizon 조합이다.
입력 pre-check가 \(\mathcal E_{\mathrm{op}}\)를 벗어나면 해당 조합은 \emph{declared OOD}이며 primary
denominator에 넣지 않는다. 범위 밖 qualitative run은 non-canonical OOD로만 기록한다.
관측된 오차나 실패를 본 뒤 in-envelope 조합을 OOD로 소급 분류하는 것은 금지한다.

반대로 pre-check를 통과한 canonical attempt에서 package rejection, current tangent rank collapse,
solver divergence/non-finite state, unresolved renderer transport 또는 covariance bound/cap 위반이 생기면
\emph{in-envelope method failure}다. 실패 frame을 valid-frame conditional 평균에 넣지는 않지만,
그 package/sequence를 canonical denominator에서 제거하지 않는다. 미리 등록한 수를 각각
\(N_{\mathrm{pkg}}^{\mathrm{attempt}}\), \(N_{\mathrm{seq}}^{\mathrm{attempt}}\)라 하고
\begin{equation}
  R_{\mathrm{accept}}
  =\frac{N_{\mathrm{pkg}}^{\mathrm{accepted}}}
         {N_{\mathrm{pkg}}^{\mathrm{attempt}}},\qquad
  R_{\mathrm{success}}
  =\frac{N_{\mathrm{seq}}^{\mathrm{successful}}}
         {N_{\mathrm{seq}}^{\mathrm{attempt}}}
  \label{eq:minimal-attempt-success-contract}
\end{equation}
을 conditional quality/timing과 함께 보고한다. Package rejection으로 실행할 수 없어진 사전 등록
sequence도 end-to-end \(N_{\mathrm{seq}}^{\mathrm{attempt}}\)에는 남고 unsuccessful로 센다.
Failure reason은 mutually exclusive first-failure code와 time-to-failure로 기록한다.
Development split에서 동결한 \(\tau_{\mathrm{fail},A}\), \(\tau_{\mathrm{fail},C}\)에 대해
\(1-R_{\mathrm{accept}}>\tau_{\mathrm{fail},A}\)이면 Gate A,
\(1-R_{\mathrm{success}}>\tau_{\mathrm{fail},C}\)이면 final deployable Gate C를 실패한다.
성공 sequence만의 낮은 평균 오차나 latency로 이 failure cap을 상쇄할 수 없다.
고정 penalty를 실패에 부여한 capped aggregate는 development에서 penalty와 sensitivity를 동결한
보조 지표로 추가할 수 있으나, 임의의 penalty 선택에 의존하므로 필수 primary가 아니며
두 rate와 conditional metric을 대체하지 않는다.
Canonical V0는 실패 시 frame별 substep 수를 바꾸거나 corrector를 자동 추가하지 않는다.

모든 canonical typed package와 run manifest는 다음 method signature를 필수 immutable field로 저장한다.
\[
\begin{aligned}
  \texttt{basis\_type} &= \texttt{generalized\_eigen\_v1},\\
  \texttt{tangent\_fit\_version} &= \texttt{tangent\_fit\_v1},\\
  \texttt{tangent\_gauge\_version} &= \texttt{attachment\_line\_v1},\\
  \texttt{bend\_projection\_version} &= \texttt{corotated\_tangent\_plane\_v2},\\
  \texttt{core\_scope} &= \texttt{flat\_rest\_v0}.
\end{aligned}
\]
다섯 field와 이들이 생성한 basis, tangent support/weight, fitted-frame projector 및
core-scope validation report는 package hash에 포함한다. Run manifest가 package의 field와 hash에
정확히 일치하지 않으면 fail-fast한다. Predicted scaffold package는 추가로
식~\eqref{eq:minimal-r1-version-contract}의 V0-R1 contract/model/decoder hash를 요구한다.
정량 core의 별도 immutable renderer protocol field는
\(\texttt{appearance\_mode=degree0\_v0}\)로 고정하고 run hash에 포함한다.
이 field가 다르거나 누락된 render는 canonical quantitative result로 사용하지 않는다.

\section{Frozen Minimal V0 architecture}

\subsection{다섯 단계}

\begin{enumerate}
  \item Static GS에서 sparse oriented anchor scaffold를 만든다.
  \item User material/attachment로 고정한 PD-style structural package와
  Global/patch-complement basis를 offline 생성한다.
  \item 매 frame 모든 anchor에서 current-surface quasi-steady wind force를 한 번 계산한다.
  \item Deterministic score로 fixed-cardinality Local unit을 선택하고
  Global+selected Local coupled reduced system을 한 번 적분한다.
  \item 최종 anchor deformation을 Gaussian mean/covariance에 affine transport하고 렌더한다.
\end{enumerate}

\begin{tcolorbox}[colback=RejectLight,colframe=RejectRed,title=V0 금지 경로,fonttitle=\bfseries]
Physics-only V0에는 patch nodal-force proposal, generalized-to-nodal lifting,
overlap force KKT/Schur, support-restricted reaction routing과 same-frame aerodynamic
corrector가 없다. Aerodynamic hyper-reduction H, missing-force network F,
learned gate G, confidence calibration network와 Teacher-FSI도 구현하지 않는다.
\end{tcolorbox}

\subsection{Causal frame dependency}

\begin{equation}
\begin{aligned}
  (q^{t-1},z^{t-1})
  &\longrightarrow (x^{t-1},v^{t-1},x_{\mathrm{pred}}^t),\\
  x_{\mathrm{pred}}^t&\longrightarrow (p^t,h^t),\\
  (x^{t-1},v^{t-1},W^t)&\longrightarrow
  (n^{t-1},a_{\mathrm{safe}}^{t-1},f_{\mathrm{aero}}^t)
  \longrightarrow\mathcal S_t^{\mathrm{solve}},\\
  (p^t,h^t,f_{\mathrm{aero}}^t,\mathcal S_t^{\mathrm{solve}})
  &\longrightarrow(\widetilde q^t,\widetilde z^t)
  \longrightarrow(q^t,z^t)
  \longrightarrow \mathcal G^t.
\end{aligned}
  \label{eq:minimal-frame-dependency}
\end{equation}

Aerodynamic force는 frame 시작 geometry에서 한 번 평가해 step 동안 hold한다.
새 geometry는 다음 frame의 force에 반영된다. 이는 의도적인 one-frame staggered coupling이며
same-frame fixed-point convergence를 주장하지 않는다.

\section{Setup: topology-distilled sparse scaffold}

\subsection{Anchor와 oriented relation}

Rendering Gaussian과 simulation anchor resolution을 분리한다.
Static Gaussian surface evidence에서 \(K\)개의 anchor \(x_i^0\)를 선택한다.
각 anchor에는 tangent frame \((t_{i1}^0,t_{i2}^0,n_i^0)\),
sparse oriented relation과 nonnegative area ownership \(a_i^0\)를 둔다.

Training에서는 source mesh topology를 privileged supervision으로 사용할 수 있지만,
target setup/runtime은 source 또는 reconstructed target triangle mesh를 입력으로 요구하지 않는다.
Network는 dense \(M,C,K\)를 직접 예측하지 않고 relation/layer score와 area ownership만 출력한다.
Sparse physical reconstruction과 surface-aligned Gaussian work
\cite{jing2026bendtwin,guedon2024sugar,yu2024gof}는 representation baseline과 경계를 정하는 데 사용한다.

\subsection{Hard-attachment 기반 in-plane line gauge}

여기서 \(n_i^0\)는 기존 static-GS thin-axis와 versioned same-layer orientation propagation이 만든
surface normal이며, 이 절은 normal을 새로 학습하지 않고 남아 있는 in-plane gauge만 고정한다.
Core hard-edge attachment anchor 집합을 \(\mathcal A_{\mathrm{att}}\)라 하고, 고정 positive weight로
중심 \(\bar x_{\mathrm{att}}\)와 covariance \(C_{\mathrm{att}}\)를 계산한다. 고윳값을
\(\lambda_1\geq\lambda_2\geq\lambda_3\geq0\)라 할 때
\begin{equation}
  C_{\mathrm{att}}g_{\mathrm{att}}=\lambda_1g_{\mathrm{att}},\qquad
  \|g_{\mathrm{att}}\|_2=1,\qquad
  L_{\mathrm{att}}=g_{\mathrm{att}}g_{\mathrm{att}}^\top,\qquad
  \frac{\lambda_1-\lambda_2}{\lambda_1+\varepsilon_{\mathrm{att}}}
  \geq\tau_{\mathrm{gap}}>0
  \label{eq:minimal-attachment-line-gauge}
\end{equation}
를 요구한다. 물리적 gauge는 vector \(g_{\mathrm{att}}\)가 아니라 sign-invariant unoriented line
\(L_{\mathrm{att}}\)다. Eigensolver 출력은 \(q_i\)와 \(L_{\mathrm{att}}\)를 계산하기 전에
unit-normalize한다. 따라서 \(L_{\mathrm{att}}\)는
\(L_{\mathrm{att}}^2=L_{\mathrm{att}}\), \(\operatorname{tr}(L_{\mathrm{att}})=1\)인
rank-one orthogonal projector이고, eigenvector의 임의 scale과 sign에 무관하다.
\(\lambda_1\) floor나 eigengap을 통과하지 못하면 world-axis에 맞춘 임의 sign 또는
secondary direction으로 우회하지 않고 package를 reject한다.

Anchor \(i\)에서는 \(P_i^0=I-n_i^0(n_i^0)^\top\),
\(q_i=P_i^0g_{\mathrm{att}}\)에 대해 \(\|q_i\|\geq\tau_{\mathrm{gauge}}\)를 요구하고
\begin{equation}
  L_i^0=\frac{q_iq_i^\top}{q_i^\top q_i},\qquad
  t_{i1}^0\in\operatorname{range}(L_i^0),\quad \|t_{i1}^0\|_2=1,\qquad
  t_{i2}^0=n_i^0\times t_{i1}^0
  \label{eq:minimal-local-line-gauge}
\end{equation}
로 둔다. \(t_{i1}^0\mapsto-t_{i1}^0\)와 \(t_{i2}^0\mapsto-t_{i2}^0\)는 같은 gauge다.
Opposite pair, feature, label, loss와 decoder는 이 동시 sign flip에 불변이어야 한다.
Byte-level serialization에 representative가 필요하면 immutable attachment-ID 순서에서 처음 만나는
nondegenerate attachment pair에만 sign을 맞추며, world-axis component sign은 사용하지 않는다.
Minimal V0의 learned in-plane correction은 \(\Delta\theta_i=0\)으로 고정하고 topology model 출력에
추가하지 않는다. Gauge algorithm, attachment IDs/weights, eigengap/floor, representative rule을
normalization/projector tolerance와 함께
\(\texttt{tangent\_gauge\_version=attachment\_line\_v1}\) package hash에 저장한다.

\subsection{고정 V0 relation schema에서 analytic constraint package로}

V0 relation schema는
\(\sigma_{\mathrm{rel}}=\texttt{edge-opposite-triplet-v1}\) 하나로 고정한다.
Accepted undirected same-layer edge \(e=(i,k)\)는 stretch relation이며
\[
  A_e^{\mathrm{str}}x=x_i-x_k,
  \qquad
  \mathcal C_e^{\mathrm{str}}
  =\{d\in\mathbb R^3:\|d\|_2=\ell_e^0\},
  \qquad
  \ell_e^0=\|x_i^0-x_k^0\|_2
\]
로 둔다. Bending은 anchor \(i\)의 각 rest tangent axis \(a\in\{1,2\}\)에서 정의한다.
Accepted neighbor 중 \(-t_{ia}^0,+t_{ia}^0\) 방향의 same-layer pair
\((k_{ia}^-,k_{ia}^+)\)를 development에서 동결한 각도 및 길이비 규칙으로 결정한다.
Boundary이거나 valid opposite pair가 없는 axis는 bending center에서 제외한다.
Retained pair는 angle/length-ratio 조건뿐 아니라 rest tangential interpolation residual cap과
해당 support의 \(\tau_{\mathrm{planar}}\) cap도 통과해야 하며,
asset별 retained bending-center coverage를 기록한다.
\(\ell_-=\|x_i^0-x_{k^-}^0\|_2\),
\(\ell_+=\|x_{k^+}^0-x_i^0\|_2\)와
\(\beta_-={\ell_+}/{(\ell_-+\ell_+)}\),
\(\beta_+={\ell_-}/{(\ell_-+\ell_+)}\)에 대해
\[
  A_{ia}^{\mathrm{bend}}x
  =x_i-\beta_-x_{k^-}-\beta_+x_{k^+},
  \qquad
  d_{ia}^0=A_{ia}^{\mathrm{bend}}x^0
\]
로 둔다. Exact collinear flat triplet에서는 \(d_{ia}^0=0\)이다.
Irregular flat sampling에서는 tangential residual cap으로 in-plane pollution을 제한하고,
flat-rest support이더라도 \(d_{ia}^0\neq0\)일 수 있음을 허용한다.
이를 임의로 0으로 만드는 flatness shortcut은 사용하지 않는다.
그 크기는
\begin{equation}
  \delta_{ia}^{\mathrm{rest}}
  =\frac{\|d_{ia}^0\|_2}
  {(\ell_-+\ell_+)/2+\varepsilon_\ell}
  \label{eq:minimal-irregular-rest-diagnostic}
\end{equation}
로 package/asset별 분포를 저장한다. 이는 irregular sampling이 남긴 rest residual을 해석하기 위한
파생 diagnostic일 뿐 별도 acceptance gate가 아니다. 이미 동결한 tangential interpolation residual,
식~\eqref{eq:minimal-flat-rest-planarity}의 planarity cap과 retained coverage가 package validity를 소유한다.
\(\delta_{ia}^{\mathrm{rest}}\)만으로 relation을 제거하거나 새 curvature/shell backend를 여는 것은
Gate A/B/C 이전 V0 범위에 포함하지 않는다.
Opposite pair가 없으면 정해진 halo를 한 번 확대하고, 그래도 실패하면 해당 boundary/axis constraint만
생성하지 않는다. Retained coverage가 development에서 동결한 minimum 아래이면
다른 constraint schema로 자동 전환하지 않고 package를 reject/user hint로 보낸다.
Stretch input은 rest-length sphere의 Euclidean closest point로 project하고, input norm이 numerical
floor 아래이면 rest edge direction을 사용한다. Bending에는 nonzero-\(d_{ia}^0\) sphere orbit을 사용하지
않는다. 그 orbit은 \((n_i^0)^\top d_{ia}^0=0\)일 때
\(d=d_{ia}^0+\epsilon n_i^0\)에 대한 minimized energy가 \(O(\epsilon^4)\)로 시작하여
normal quadratic stiffness를 소실시키기 때문이다.

Bending의 유일한 projector는 다음처럼 고정한다.
\[
\begin{gathered}
  \texttt{bend\_projection\_version}\\[-0.25em]
  =\texttt{corotated\_tangent\_plane\_v2}.
\end{gathered}
\]
식~\eqref{eq:minimal-surface-tangent-fit}의 고정 support, weight와 rest coordinate를 재사용해
\(x^0\)와 \(x_{\mathrm{pred}}^t\)에서 각각 \(J_i^0,J_i^{t,\mathrm{pred}}\)를 계산한다. 두 상태 모두
동일한 orthonormalization을 적용한다.
\begin{equation}
\begin{aligned}
  \mathcal R_i(J)
  &=\left[
    \frac{J(:,1)}{\|J(:,1)\|_2},\;
    n_i(J)\times\frac{J(:,1)}{\|J(:,1)\|_2},\;
    n_i(J)
  \right],\\
  n_i(J)
  &=s_i\frac{J(:,1)\times J(:,2)}{\|J(:,1)\times J(:,2)\|_2},\\
  R_i^0&=\mathcal R_i(J_i^0),\qquad
  R_i^{t,\mathrm{pred}}=\mathcal R_i(J_i^{t,\mathrm{pred}}),\qquad
  Q_i^t=R_i^{t,\mathrm{pred}}(R_i^0)^\top .
\end{aligned}
  \label{eq:minimal-bending-corotated-frame}
\end{equation}
Stored \((t_{i1}^0,t_{i2}^0,n_i^0)\)와 fitted rest frame이 정확히 같다고 가정하지 않는다.
Rest와 predictor에 같은 construction을 사용하므로 rest에서 \(Q_i^0=I\)이고 proper rigid transform에서
\(R_i^{t,\mathrm{pred}}=QR_i^0\)다. 기존 tangent-fit rank/condition/normal floor를 그대로 적용하며
실패를 임의 frame fallback으로 우회하지 않는다.

Fitted normal을
\(n_{i,\mathrm{fit}}^0=R_i^0(:,3)\),
\(n_i^{t,\mathrm{pred}}=R_i^{t,\mathrm{pred}}(:,3)\)라 하고
\begin{equation}
\begin{aligned}
  b_{ia}^0
  &=\bigl((n_{i,\mathrm{fit}}^0)^\top d_{ia}^0\bigr)n_{i,\mathrm{fit}}^0,
  &b_{ia}^t&=Q_i^t b_{ia}^0,\\
  \mathcal C_{ia}^{\mathrm{bend},t}
  &=\{b_{ia}^t+v:(n_i^{t,\mathrm{pred}})^\top v=0\},
  &\Pi_{ia}^{\mathrm{bend},t}(d)
  &=d-n_i^{t,\mathrm{pred}}(n_i^{t,\mathrm{pred}})^\top(d-b_{ia}^t).
\end{aligned}
  \label{eq:minimal-bending-corotated-projector}
\end{equation}
로 target을 한 번 계산해 해당 frame의 global step 동안 freeze한다. Unit fixture는 frame을 \(R_i^0\)로
고정한 뒤 \(d_{ia}^0\pm\epsilon n_{i,\mathrm{fit}}^0\)의 residual이 \(\pm\epsilon n_{i,\mathrm{fit}}^0\)이고
central finite-difference stiffness가 \(w_{ia}^{\mathrm{bend}}\)인지 검사한다. Integrated fixture는 frame을
다시 fit하여 rest exactness와 random proper-rigid reproduction을 검사하고, normal perturbation에는
positive non-collapsed response를 요구한다. Integrated stiffness를 정확히
\(w_{ia}^{\mathrm{bend}}\)로 요구하려면 추가로 \(\delta J_i=0\)인 perturbation이어야 한다.
Relation support, \(A_{ia}^{\mathrm{bend}}\), owner area, \(w_{ia}^{\mathrm{bend}}\)와 Hessian은 고정이며
새 local/global iteration이나 runtime stage를 만들지 않는다. Projector version, fitted-frame convention,
threshold, test tolerance와 replay hash를 structural package에 저장한다.

Anchor area는 constraint weight에 중복 없이 들어가도록 nonnegative owner share로 분배한다.
Accepted stretch degree를 \(d_i^{\mathrm{str}}\), anchor \(i\)의 retained bending-axis 수를
\(n_i^{\mathrm{bend}}\)라 두고
\[
  a_e^{\mathrm{own}}
  =\frac{a_i^0}{d_i^{\mathrm{str}}}+\frac{a_k^0}{d_k^{\mathrm{str}}},
  \qquad
  a_{ia}^{\mathrm{own}}=\frac{a_i^0}{n_i^{\mathrm{bend}}}
\]
로 고정한다. Stretch share의 합은 valid anchor area total이며,
bending share의 합은 retained bending-center area total이다.
Geometric boundary에서 제외한 bending area를 interior constraint에 재분배하지 않고
coverage ratio로 별도 보고한다. Required interior anchor에서 어느 분모가 0이면 package를 accept하지 않는다.
Stretch edge와 bending triplet의 owner area를 각각
\(a_e^{\mathrm{own}}\), \(a_{ia}^{\mathrm{own}}\),
\(\bar\ell_{ia}=(\ell_-+\ell_+)/2\)라 하면
\[
  w_e^{\mathrm{str}}
  =\kappa_{\mathrm{str}}^{(\mathsf u)}
   \frac{a_e^{\mathrm{own}}}{(\ell_e^0)^2+\varepsilon_\ell^2},
  \qquad
  w_{ia}^{\mathrm{bend}}
  =\kappa_{\mathrm{bend}}^{(\mathsf u)}
   \frac{a_{ia}^{\mathrm{own}}}{\bar\ell_{ia}^{\,4}+\varepsilon_\ell^4}.
\]
SI에서 \(\kappa_{\mathrm{str}}\)는 surface stretching stiffness \([\mathrm{N}/\mathrm m]\),
\(\kappa_{\mathrm{bend}}\)는 bending rigidity \([\mathrm{N}\,\mathrm m]\)이고,
nondimensional branch는 같은 식의 dimensionless preset을 사용한다.
따라서 두 channel의 \(w_j\)는 모두 quadratic position residual에 맞는 \([\mathrm N/\mathrm m]\) 또는
그 nondimensional 대응량이다.

이 고정 변환을 요약하면

\begin{equation}
\begin{aligned}
  \mathcal P_{\mathrm{V0}}
  &=\left(
    \mathcal R_{\mathrm{str}},\mathcal R_{\mathrm{bend}},
    \{\operatorname{supp}(A_j),A_j,\Pi_j,w_j\}_{j\in\mathcal R}
  \right),\\
  \mathcal P_{\mathrm{V0}}
  &=
  \mathcal D_{\texttt{edge-opposite-triplet-v1}}
  \left(\mathcal R_{\mathrm{acc}},T^0,a^0,
  \vartheta_{\mathrm{mat}},\mathsf u\right),\\
  \mathcal R&=\mathcal R_{\mathrm{str}}\cup\mathcal R_{\mathrm{bend}}.
\end{aligned}
  \label{eq:minimal-topology-to-constraints}
\end{equation}

으로 고정한다. Support ordering, neighbor-selection threshold, owner share,
projection-set version과 stiffness preset을 package에 저장한다.
Network가 constraint, stiffness 또는 dense operator를 예측하는 것은 아니다.
Alternative hinge/hyperedge schema는 이 V0의 per-asset 선택지가 아니라 후속 backend version이다.

\subsection{Minimal validity check}

\begin{equation}
\begin{aligned}
  &a_i^0>0,\qquad
  \sum_i a_i^0 \approx A_{\mathrm{asset}},\qquad
  m_i=\rho_A a_i^0,\qquad
  \sum_i m_i \approx M_{\mathrm{asset}},\\
  &M=\operatorname{blkdiag}_{i\in V}(m_i I_3),\qquad
  \texttt{mass\_matrix\_type}=\texttt{lumped\_anchor\_v1}.
\end{aligned}
  \label{eq:minimal-area-mass}
\end{equation}

Structural package hash는 \texttt{mass\_matrix\_type}, stable anchor ordering에 따른 ordered
\(m_i\) 또는 동등한 repeated diagonal, hard-attachment free map을 소유한다. 이 정보로 \(M\)과
\(M_f=N_f^\top MN_f\)를 정확히 재현할 수 있으면 충분하며 dense \(M\) array 저장을 요구하지 않는다.

Canonical Minimal V0의 area total은 별도 calibrated surface-area estimator가 아니라
\[
  L_{\mathrm{ref}}=\operatorname{diameter}(x^0)>0,\qquad
  A_{\mathrm{ref}}=L_{\mathrm{ref}}^2,\qquad
  A_{\mathrm{asset}}:=A_{\mathrm{ref}},\qquad
  M_{\mathrm{asset}}:=\rho_AA_{\mathrm{ref}}
\]
인 deterministic package convention으로 고정한다. 이는 physical absolute surface-area 또는
engineering mass 측정값이라는 주장이 아니며, 목적은 Gaussian/anchor density가 달라도 relative
ownership과 motion scale이 크게 바뀌지 않게 하는 것이다. Object-normalized nondimensional geometry에서
\(L_{\mathrm{ref}}=1\)이면 \(A_{\mathrm{ref}}=1\)이다. 공통 area scale에 대한 trajectory invariance는
그 scale을 소유하는 \(M,C,K,\bar h,f_{\mathrm{aero}}\)가 함께 similarity-scale되고
관련 floor, epsilon과 clamp가 비활성이거나 같은 typed rule로 scale되는 조건에서만 요구한다.
Nondimensional branch의 \(a_i^0,\rho_A,m_i,A_{\mathrm{asset}},M_{\mathrm{asset}}\)는
\((L_0,M_0,T_0)\)로 정규화한 dimensionless package 값이다.

V0 validity는 다음 hard check만 사용한다.

\begin{itemize}
  \item local relation rank/condition과 tangent orientation
  \item close front/back layer를 잇는 shortcut의 부재
  \item positive area/mass와 total ownership
  \item hard attachment가 유효 anchor 집합에 놓였는지
\end{itemize}

Relation score threshold는 development split에서 한 번 동결한다.
NLL/Brier/ECE/AUPRC calibration network를 만들지 않으며,
invalid package는 reject, user hint 또는 explicit scaffold fallback으로 보낸다.

\subsection{단순 structural package}

V0 backend는 \emph{one-local-step fixed-Hessian PD-style stretch+bending graph} 하나로 고정한다
\cite{bouaziz2014projective}.
Typed package는 backend/schema version, \(\mathcal R_{\mathrm{str}}\),
\(\mathcal R_{\mathrm{bend}}\), 각 \(\operatorname{supp}(A_j),A_j,\Pi_j,w_j\),
free-DOF map, damping preset과 허용 timestep을 저장한다.
Stretch의 fixed sphere와 bending의 predictor-frozen corotated tangent-plane target에 대한
versioned projector가 유일한 local step이다. \(A_j,w_j\)와
\(K=\sum_jw_jA_j^\top A_j\)는 setup 뒤 고정된다.
별도 corotational graph, nonlinear tangent update 또는 alternative bending projector를
병렬 mainline으로 두지 않는다.
별도 thin-shell novelty를 주장하지 않는다.

Frame 시작에
\[
  x_{\mathrm{pred}}^t=x^{t-1}+\Delta t\,v^{t-1},\qquad
  p_j^t=\Pi_j^t(A_jx_{\mathrm{pred}}^t)
\]
로 projector target을 한 번 계산하고 hold한다.
Global step의 quadratic energy는

\begin{equation}
  E_{\mathrm{str}}^t(x)
  =
  \frac12\sum_{j\in\mathcal R}
  w_j\left\|A_jx-p_j^t\right\|_2^2,
  \qquad
  K=\sum_j w_jA_j^\top A_j,\quad
  h^t=\sum_j w_jA_j^\top p_j^t,\quad
  \bar h^t=h^t-Kx^0.
  \label{eq:minimal-structural-backend}
\end{equation}

Runtime reduced coordinate는 rest anchor에서의 displacement를 나타내므로
아래 coupled step은 raw target \(h^t\)가 아니라 rest offset을 제거한
\(\bar h^t\)를 소비한다. 따라서 zero-wind rest state에서 불필요한 constant load가 생기지 않는다.

Hard attachment는 setup에서 free-DOF elimination 또는 attachment를 만족하는 basis로 처리한다.
Runtime support reaction을 복원하지 않는다.
Damping은 하나의 고정 Rayleigh/modal preset
\(C=\alpha M+\beta K\)로 둔다. Material UI는 stiffness scale과 damping preset의
단조 제어만 보장하면 충분하다.

Structural hard acceptance는 다음 네 가지다.

\begin{itemize}
  \item rigid transform sanity에서 눈에 띄는 가짜 restoring artifact가 없음
  \item zero wind에서 rest 쪽으로 안정적으로 감쇠
  \item target visual band의 dominant sway/flutter가 완전히 소실되지 않음
  \item 10--30초 rollout에서 NaN, scale explosion 또는 발산이 없음
\end{itemize}

Exact energy/action tangent, rest-linear equality, exact support wrench와
8--16 mode 정밀 일치는 diagnostic일 수 있으나 package acceptance의 hard gate가 아니다.

\section{Offline Global/Local basis}

\subsection{Global basis}

V0 canonical Global basis는
\(\texttt{basis\_type=generalized\_eigen\_v1}\) 하나로 고정한다.
Free-coordinate fixed-Hessian package operator의 generalized eigenproblem에서 eigenvalue가 낮은
accepted \(r_g\)개 mode를 순서대로 사용한다 \cite{barbic2005stvk,diener2009windbasis}.
Hard-attachment free map을 \(N_f\)라 두면
\(M_f=N_f^\top MN_f\), \(K_f=N_f^\top KN_f\)이고,
full anchor basis는 \(\Phi=N_f\widetilde\Phi\)로 lift한다.

\begin{equation}
  K_{\mathrm f}\widetilde\Phi
  =
  M_{\mathrm f}\widetilde\Phi\Lambda,\qquad
  \Phi^\top M\Phi=I.
  \label{eq:minimal-global-basis}
\end{equation}

여기서 \(K_f\)는 predictor-frozen one-local-step PD package가 저장한 fixed Hessian이다.
Minimized nonlinear constraint energy의 exact rest tangent라고 주장하지 않으며, 이 근사를 이유로
tangent-consistent basis, 새 shell backend 또는 추가 nonlinear iteration을 V0에 넣지 않는다.

Global은 전체 sway, large bending과 low-frequency response를 소유한다.
Basis ordering, repeated-eigenvalue tie-break와 column sign convention을 development에서 고정하고
basis array와 함께 package hash에 포함한다. POD는 Gate A/B/C 이후의 optional ablation일 뿐이며,
canonical V0 result나 matched-backend baseline의 basis를 대체할 수 없다.
Strong Global rank sweep은 같은 \texttt{generalized\_eigen\_v1} ordering을 prefix로 사용하며
MC2의 필수 baseline이다.

\subsection{Patch-local complementary basis}

Setup에서 고정한 core--halo patch \(r\)의 free-anchor-DOF injection을
\(E_r\in\mathbb R^{3K\times n_r}\)라 한다.
Global complement를 만든 뒤 dense tail을 허용하지 않고,
patch support 안에서 orthogonality constraint를 만족하는 restricted mode를 푼다.
\(M_{rr}=E_r^\top ME_r\), \(K_{rr}=E_r^\top KE_r\)와
\(C_r=\Phi^\top ME_r\)를 만들고, \(Z_r\)를 \(C_rZ_r=0\)의 orthonormal null basis라 하면

\begin{equation}
\begin{aligned}
  (Z_r^\top K_{rr}Z_r)V_r
  &=
  (Z_r^\top M_{rr}Z_r)V_r\Lambda_r,\\
  \Psi_r&=E_rZ_rV_r,\qquad
  V_r^\top Z_r^\top M_{rr}Z_rV_r=I.
\end{aligned}
  \label{eq:minimal-local-complement}
\end{equation}

\(V_r\)는 lowest accepted \(k_r\) columns로 고정하고, \(k_r\), halo와 pruning rule을
development package에서 동결한다. 따라서 Gate B에서 asset별로 Local candidate를 cherry-pick하지 않는다.

\begin{equation}
  \Phi^\top M\Psi_r\approx0,\qquad
  \Psi_r^\top M\Psi_r=I.
  \label{eq:minimal-complement-check}
\end{equation}

서로 다른 patch basis를 전역 QR로 섞지 않는다.
Patch끼리는 강제 직교화하지 않고 \(M_{rs},C_{rs},K_{rs}\) cross block을 저장한다.
Null dimension이 부족하거나 strong overlap이 rank deficiency를 만들면
setup에서 halo 확대, patch 병합 또는 column pruning 중 하나를 기록하고 재검증한다.
이는 adaptive subspace 및 complementary dynamics
\cite{benchekroun2023complementary,xiang2026freeform}와 비교한다.

\subsection{Persistent state}

모든 \(\Psi_r\)는 fixed coordinate slot을 가진다.
Runtime은 \(z_{\mathrm{sup}},\dot z_{\mathrm{sup}}\)를 유지하며,
active/decay subset만 gather--solve--scatter한다.
Basis를 frame마다 다시 만들거나 state를 remap하지 않는다.

Anchor displacement와 velocity의 유일한 reconstruction은

\begin{equation}
  u^t=\Phi q^t+\sum_r\Psi_rz_r^t,\qquad
  x^t=x^0+u^t,\qquad
  v^t=\Phi\dot q^t+\sum_r\Psi_r\dot z_r^t.
  \label{eq:minimal-state-reconstruction}
\end{equation}

Inactive slot은 정확히 zero state다. Decay slot은 아래 release tolerance를 통과할 때까지
reconstruction과 coupled solve에 남으며, 통과한 뒤에만 \((z_r,\dot z_r)\)를 zero로 만들고 제외한다.

\section{Runtime surface와 full-anchor aerodynamics}

\subsection{Current surface map}

Anchor \(i\)는 setup에서 고정한 ordered same-layer support
\(\mathcal N_i^{\mathrm{surf}}\)와 positive normalized weight
\(\omega_{ik}^{\mathrm{surf}}>0\),
\(\sum_{k\in\mathcal N_i^{\mathrm{surf}}}\omega_{ik}^{\mathrm{surf}}=1\)을 가진다.
Rest와 current에 서로 다른 centering을 쓰되 rest 2D coordinate는 고정한다.

\begin{equation}
\begin{aligned}
  \bar x_i^s
  &=\sum_{k\in\mathcal N_i^{\mathrm{surf}}}
    \omega_{ik}^{\mathrm{surf}}x_k^s,
    \qquad s\in\{0,t\},\\
  \xi_{ik}^{\mathrm{surf}}
  &=(T_i^0)^\top(x_k^0-\bar x_i^0),
  \qquad
  \sum_k\omega_{ik}^{\mathrm{surf}}\xi_{ik}^{\mathrm{surf}}=0,\\
  G_{i,\mathrm{surf}}^0
  &=\sum_k\omega_{ik}^{\mathrm{surf}}
    \xi_{ik}^{\mathrm{surf}}(\xi_{ik}^{\mathrm{surf}})^\top,\\
  C_{i,\mathrm{surf}}^s
  &=\sum_k\omega_{ik}^{\mathrm{surf}}
    (x_k^s-\bar x_i^s)(\xi_{ik}^{\mathrm{surf}})^\top,\\
  J_i^s
  &=C_{i,\mathrm{surf}}^s(G_{i,\mathrm{surf}}^0)^{-1},\\
  J_i^s
  &=\arg\min_{J\in\mathbb R^{3\times2}}
    \sum_k\omega_{ik}^{\mathrm{surf}}
    \left\|(x_k^s-\bar x_i^s)-J\xi_{ik}^{\mathrm{surf}}\right\|_2^2.
\end{aligned}
  \label{eq:minimal-surface-tangent-fit}
\end{equation}

이 positive-weight, weighted-centering, rest-coordinate least-squares convention을
\(\texttt{tangent\_fit\_v1}\)로 고정하며, 식~\eqref{eq:minimal-affine-weights}--
\eqref{eq:minimal-affine-fit}의 Gaussian binding도 같은 convention을 재사용한다.
\(J\in\mathbb R^{3\times2}\)의 singular value를
\(\sigma_1(J)\geq\sigma_2(J)\geq0\),
\(G\in\mathbb R^{2\times2}\)의 eigenvalue를
\(\lambda_1(G)\geq\lambda_2(G)\geq0\)로 정렬하고, shared numerical-rank contract를
\begin{equation}
\begin{aligned}
  \sigma_2(J)
  &\geq\max\!\left(\sigma_{J,\mathrm{abs}},\tau_J\sigma_1(J)\right),
  &\sigma_{J,\mathrm{abs}}&>0,\quad 0<\tau_J<1,\\
  \lambda_2(G)
  &\geq\max\!\left(\lambda_{G,\mathrm{abs}}^{(\mathsf u)},
                    \tau_G\lambda_1(G)\right),
  &\lambda_{G,\mathrm{abs}}^{(\mathsf u)}&>0,\quad 0<\tau_G<1
\end{aligned}
  \label{eq:minimal-tangent-numerical-rank}
\end{equation}
로 고정한다. \(\sigma_{J,\mathrm{abs}}\)는 dimensionless이고,
\(\lambda_{G,\mathrm{abs}}^{(\mathsf u)}\)는 SI에서 \(L^2\), nondimensional branch에서
dimensionless인 typed floor이며 SI-to-nd 변환은
\(\widehat\lambda_{G,\mathrm{abs}}=\lambda_{G,\mathrm{abs}}/L_0^2\)다.
각 consumer의 condition cap과 이 threshold를 함께 동결한다. 모든 \(G^{-1}\)은 \(G\)가 이 식과
condition cap을 통과한 뒤에만 계산하고, \(((J^0)^\top J^0)^{-1}\), normal과 area consumer는
해당 \(J\)가 같은 식을 통과한 뒤에만 계산한다. Threshold, unit branch와 CPU/GPU replay 결과를
\(\texttt{tangent\_fit\_v1}\) package hash에 저장한다.
Flat-rest core의 planarity statistic은
\begin{equation}
  \eta_{i,\mathrm{planar}}^0
  =
  \frac{
    \max_{k\in\mathcal N_i^{\mathrm{surf}}}
    \left|(x_k^0-\bar x_i^0)^\top n_i^0\right|
  }{
    \max_{k\in\mathcal N_i^{\mathrm{surf}}}
    \|x_k^0-\bar x_i^0\|_2
  }
  \leq\tau_{\mathrm{planar}},
  \label{eq:minimal-flat-rest-planarity}
\end{equation}
로 고정한다. Denominator가 0인 support는 아래 rank check보다 먼저 package reject한다.
Ordered tangent column의 orientation sign \(s_i\in\{-1,+1\}\)는 rest fitted normal이
stored \(n_i^0\)와 같은 방향이 되도록 setup에서 고정한다. 그러면

\begin{equation}
\begin{aligned}
  \widetilde n_i^s
  &=s_i\bigl(J_i^s(:,1)\times J_i^s(:,2)\bigr),\qquad s\in\{0,t\},\\
  \rho_{a,i}^{s,\mathrm{raw}}
  &=\frac{\|\widetilde n_i^s\|}{\|\widetilde n_i^0\|},\qquad
  \bar\rho_{a,i}^s
  =\operatorname{clip}\!\left(
    \rho_{a,i}^{s,\mathrm{raw}};
    \rho_{a,\min},\rho_{a,\max}\right),\\
  a_{i,\mathrm{safe}}^s
  &=a_i^0\bar\rho_{a,i}^s,\qquad
  n_i^s=\frac{\widetilde n_i^s}{\|\widetilde n_i^s\|},
  \qquad s\in\{0,t\}.
\end{aligned}
  \label{eq:minimal-area-normal}
\end{equation}

Setup에서는 \(G_{i,\mathrm{surf}}^0\)와 \(J_i^0\)가
식~\eqref{eq:minimal-tangent-numerical-rank}을 통과하고 다음 조건을 만족해야 한다.
\[
\begin{aligned}
  \kappa(G_{i,\mathrm{surf}}^0)&\leq\kappa_{G,\mathrm{surf}},&
  \kappa((J_i^0)^\top J_i^0)&\leq\kappa_{J,\mathrm{surf}},\\
  \|\widetilde n_i^0\|&>\tau_{n,\mathrm{surf}}>0,
  &\text{orientation agreement}&=\text{pass}.
\end{aligned}
\]
Runtime에는 \(J_i^t\)가 같은 numerical-rank 식을 통과하고
\[
  \kappa((J_i^t)^\top J_i^t)\leq\kappa_{J,\mathrm{surf}},\qquad
  \|\widetilde n_i^t\|>\tau_{n,\mathrm{surf}}
\]
를 만족해야 한다.
Rest check 실패는 package reject, current check 실패는 canonical frame fail-fast다.
Surface/aero consumer는 arbitrary closest-frame 또는 rigid fallback으로 load geometry를 만들지 않는다.
그 fallback은 퇴화한 normal/area에서 존재하지 않은 aerodynamic load를 만들어 canonical quality/timing
sample을 오염시키므로, renderer transport의 visual fallback과 동일하게 취급하지 않는다.
Surface package에는 ordered \(\mathcal N_i^{\mathrm{surf}}\), \(\omega_{ik}^{\mathrm{surf}}\),
\(T_i^0\), \(\xi_{ik}^{\mathrm{surf}}\), \(G_{i,\mathrm{surf}}^0\), \(s_i\)와 모든 cap을 저장한다.
Current failure reason과 rate는 trace에 기록한다. In-envelope invalid frame은 valid-frame conditional
metric에는 넣지 않지만 해당 sequence는 식~\eqref{eq:minimal-attempt-success-contract}의
denominator에 method failure로 남긴다.
Numerical epsilon을 분모에 더해 unit normal이나 rest-area identity를 조용히 깨뜨리지 않는다.
Operating envelope의 \(\mathcal B_{\mathrm{area/affine/cov}}\) 중 area bound가
\(0<\rho_{a,\min}\leq1\leq\rho_{a,\max}<\infty\)를 소유한다.
Valid frame에서 \(\rho_{a,i}^{s,\mathrm{raw}}\)만 이 범위로 clamp하며,
raw ratio는 diagnostic과 clamp-event 판정에만, \(a_{i,\mathrm{safe}}^s\)는 aerodynamic consumer에만 사용한다.
Event \(\bar\rho_{a,i}^s\neq\rho_{a,i}^{s,\mathrm{raw}}\)는 run 전에 등록한 required
surface/aero anchor--frame evaluation 전체를 고정 분모로 하는 기존 generic area clamp-rate trace와 cap에 포함한다.
Required evaluation의 current-surface check가 실패하면 기존 method failure로 남기며 분모에서 삭제하지 않는다.
새 threshold, index 체계 또는 gate를 만들지 않는다.
Rank-2 tangent map에 3D cofactor를 적용하지 않는다.
식~\eqref{eq:minimal-surface-tangent-fit}의 \(s\)는 generic state index다.
Aerodynamic step \(t\)에서는 current index에 \(s=t-1\)을 대입해
\((J_i^{t-1},n_i^{t-1},a_{i,\mathrm{safe}}^{t-1})\)만 사용하며
predicted/new \(x^t\)로 다시 fit하지 않는다.

\subsection{Quasi-steady wind force}

Relative wind와 tangential component를

\begin{equation}
  v_{\mathrm{rel},i}^t=W(x_i^{t-1},t)-v_i^{t-1},\qquad
  v_{n,i}^t=(v_{\mathrm{rel},i}^t)^\top n_i^{t-1},\qquad
  v_{\tau,i}^t=v_{\mathrm{rel},i}^t-v_{n,i}^tn_i^{t-1}
\end{equation}

로 둔다. V0 analytic law는

\begin{equation}
\begin{aligned}
  f_{\mathrm{raw},i}^t
  &=
  \frac12\rho_{\mathrm{air}}a_{i,\mathrm{safe}}^{t-1}
  \left[
    C_n|v_{n,i}^t|v_{n,i}^tn_i^{t-1}
    +C_\tau\|v_{\tau,i}^t\|v_{\tau,i}^t
  \right],\\
  f_{\mathrm{aero},i}^t
  &=
  \begin{cases}
    f_{\mathrm{raw},i}^t,
      & \|f_{\mathrm{raw},i}^t\|_2\leq f_{\max},\\
    f_{\max}\dfrac{f_{\mathrm{raw},i}^t}{\|f_{\mathrm{raw},i}^t\|_2},
      & \|f_{\mathrm{raw},i}^t\|_2>f_{\max}.
  \end{cases}
\end{aligned}
  \label{eq:minimal-aero}
\end{equation}

모든 anchor에서 raw force를 계산하며 hyper-reduction하지 않는다. Clamp 전에
\(f_{\mathrm{raw},i}^t\)가 finite인지 검사하고, in-envelope non-finite raw force는 대체값이나 clamp로
숨기지 않는 method failure다. 위 Euclidean radial hard clamp만 force clamp로 사용하며,
\(\|f_{\mathrm{raw},i}^t\|_2>f_{\max}\)인 각 event를 식~\eqref{eq:minimal-operating-envelope}에서 이미
동결한 generic force clamp-rate trace와 cap에 포함한다. Run 전에 required raw-force anchor--frame
evaluation index set \(\mathcal I_f\)를 등록하고
\[
  r_{f,\mathrm{clip}}
  =
  \frac{1}{|\mathcal I_f|}
  \sum_{(i,t)\in\mathcal I_f}
  \mathbf 1\!\left[
    f_{\mathrm{raw},i}^t\ \text{finite}
    \ \land\ \|f_{\mathrm{raw},i}^t\|_2>f_{\max}
  \right]
\]
로 기존 force clamp-rate cap을 집계한다. \(\mathcal I_f\)의 어느 required evaluation이라도 non-finite이면
해당 attempt는 별도 method failure이며, 그 index만 분모에서 삭제해 successful finite-only run으로
재분류하지 않는다. 별도 threshold나 gate는 만들지 않는다.
Zero relative wind에서 force가 0이고, normal flip/two-sided convention과
finite-force clamp가 unit test를 통과해야 한다.
정확한 lift/wake force를 주장하지 않는다.

\section{Deterministic selection과 최소 state policy}

\subsection{Atomic Local unit}

Patch 하나 또는 strong-overlap patch를 setup에서 병합한 group을 atomic unit
\(\mathcal G\)로 둔다. Unit grouping은 runtime에 바꾸지 않는다.
Setup package는 stable unit ID/order, 각 unit의 stable ordered member list,
versioned merge/pruning rule과 그 결과로 retained된 output column range/order를
다음 네 metadata field로 저장한다.
\begin{itemize}
  \item \texttt{atomic\_group\_map}
  \item \texttt{stable\_unit\_order}
  \item \texttt{group\_member\_order}
  \item \texttt{group\_column\_ranges}
\end{itemize}
\(\Psi_{\mathcal G}\)는 이 mapping이 가리키는 serialized retained group-basis array 자체이고,
그 column order와 persistent coordinate slot은 package build 뒤 바뀌지 않는다.
Versioned setup rule은 retained member basis의 조합 또는 merged-patch basis 재생성 중 선택한 한 규칙을
명시할 수 있으나, 모든 member basis를 무조건 concatenate하도록 강제하지 않는다.
Basis array, column provenance/range와 persistent slot mapping이 다르면 다른 package hash다.

이미 계산한 dynamic anchor load를 사용해

\begin{equation}
  g_{\mathcal G}^t
  =
  \Psi_{\mathcal G}^\top
  f_{\mathrm{aero}}^t.
  \label{eq:minimal-local-drive}
\end{equation}

\(m_{\mathcal G}^{\mathrm{own}}\)은 halo를 제외한 atomic unit core anchor mass의 합으로 고정한다.
Canonical score는 cached isolated external-load compliance proxy다.

\begin{equation}
\begin{aligned}
  A_{\mathcal G}^{\mathrm{score}}
  &=
  \frac{4}{\Delta t^2}M_{\mathcal G\mathcal G}
  +\frac{2}{\Delta t}C_{\mathcal G\mathcal G}
  +K_{\mathcal G\mathcal G},\\
  \eta_{\mathcal G}^t
  &=
  \left[
  \frac{(g_{\mathcal G}^t)^\top
  (A_{\mathcal G}^{\mathrm{score}})^{-1}g_{\mathcal G}^t}
  {m_{\mathcal G}^{\mathrm{own}}+\varepsilon_m}
  \right]^{1/2}.
\end{aligned}
  \label{eq:minimal-selector}
\end{equation}

V0는 dwell-reserved capacity-\(K_A\) TopK만 사용한다.
아직 minimum dwell 안인 unit을 먼저 유지하고 남은 slot만 \(\eta\)로 채운 결과를
desired set \(\mathcal A_t^{\mathrm{des}}\)라 한다.
\(K_A\)는 fixed active cardinality이지 fixed compute budget이 아니므로,
active/solve rank와 p95 latency를 함께 보고한다.
\(K_A\), stable unit order와 고정 tie-break는 selection-config subhash 및 run hash에 포함하며,
package의 atomic-group mapping/hash와 일치하지 않으면 fail-fast한다.
Curvature 또는 random selector는 baseline일 뿐 feature fusion mainline이 아니다.
이 proxy는 structural target \(\bar h^t\), Global/inter-group coupling과 future error를 보지 않는
의도적인 근사이며, teacher-ranked gap과 activation delay로만 유효성을 판정한다.

\subsection{Active/Decay 두 상태}

\begin{itemize}
  \item \textbf{Active}: 아래 admission을 통과한 solve 직전 실제 집합
  \(\mathcal A_t^{\mathrm{pre}}\)이며
  direct projected excitation을 받는다.
  \item \textbf{Decay}: solve 직전 집합을 \(\mathcal D_t^{\mathrm{pre}}\)라 하며,
  새 Local excitation은 fade-out하지만 coupled physical solve와 damping은 유지한다.
\end{itemize}

Desired set과 실제 Active set은 다음 순서로 연결한다.

\begin{enumerate}
  \item \(\mathcal A_t^{\mathrm{des}}\cap\mathcal D_{t-1}^{\mathrm{end}}\)를 score/tie-break 순서로
  state reset 없이 먼저 재활성화한다. Active가 가득 차 있으면
  \(\mathcal A_t^{\mathrm{des}}\)에 없는 최저-priority Active unit을 골라,
  reactivated unit이 비운 Decay slot으로 보내는 atomic swap을 수행한다.
  \item 비어 있는 Active slot에는 새 desired unit을 바로 admit한다.
  \item Active replacement가 필요하면 outgoing unit을 Decay로 옮길 빈 slot이 있을 때만
  새 desired unit을 admit한다.
  \item Decay capacity가 가득 차면 outgoing Active를 그대로 유지하고 새 candidate의
  activation delay를 기록한다.
\end{enumerate}

따라서 충분한 valid candidate가 있을 때 \(|\mathcal A_t^{\mathrm{pre}}|=K_A\)지만,
capacity pressure에서는 \(\mathcal A_t^{\mathrm{pre}}\neq\mathcal A_t^{\mathrm{des}}\)일 수 있다.
이는 hidden adaptive budget이 아니라 energetic reset을 막는 deterministic admission policy다.

Deterministic parameter
\[
  \vartheta_{\mathrm{state}}
  =
  \left(K_A,K_D,N_{\mathrm{dwell}},N_{\mathrm{fade}},N_{\mathrm{hold}},
  \tau_E,\tau_x,\tau_v,\tau_{\mathrm{pop}}\right)
\]
를 development split에서 동결하고 package/run manifest에 저장한다.
\(K_A,K_D\)와 세 frame-count parameter는 모두 1 이상인 정수이고,
네 release threshold \(\tau_E,\tau_x,\tau_v,\tau_{\mathrm{pop}}\)는 모두 finite positive다.
Strict \(<\) release predicate와 0 threshold를 함께 허용하지 않으며 config validation에서 fail-fast한다.
Reference는
\[
\begin{aligned}
  M_{r,\mathrm{ref}}&=m_r^{\mathrm{own}}>0,&
  L_{\mathrm{ref}}&=\operatorname{diameter}(x^0)>0,\\
  V_{\mathrm{ref}}&=\max(W_{\max},V_{\mathrm{floor}})>0,&
  E_{r,\mathrm{ref}}&=M_{r,\mathrm{ref}}V_{\mathrm{ref}}^2.
\end{aligned}
\]
로 package에서 고정한다.
아직 \(N_{\mathrm{dwell}}\)을 채우지 않은 Active unit을 먼저 \(K_A\) slot에 유지하고,
나머지 slot을 score와 고정 tie-break 순서로 채운다.
Unit \(r\)가 마지막으로 Active였던 frame을 \(t_{r,\mathrm{off}}\)라 한다.
따라서 첫 Decay frame은 \(t=t_{r,\mathrm{off}}+1\)이고 canonical fade는

\begin{equation}
  \gamma_r^t
  =
  \max\!\left(0,1-\frac{t-t_{r,\mathrm{off}}}{N_{\mathrm{fade}}}\right),
  \qquad
  \widehat E_r^t
  =
  \frac{
    \tfrac12(\dot{\widetilde z}_r^{\,t})^\top
      M_{rr}^{\mathrm{loc}}\dot{\widetilde z}_r^{\,t}
    +\tfrac12(\widetilde z_r^t)^\top K_{rr}^{\mathrm{loc}}\widetilde z_r^t
  }{E_{r,\mathrm{ref}}+\varepsilon_E},
  \quad
  \begin{aligned}
  M_{rr}^{\mathrm{loc}}&=\Psi_r^\top M\Psi_r,\\
  K_{rr}^{\mathrm{loc}}&=\Psi_r^\top K\Psi_r .
  \end{aligned}
  \label{eq:minimal-decay-release}
\end{equation}

여기서 \(\widehat E_r^t\)는 inter-patch/Global cross term을 소유하지 않는
scale-aware \emph{release proxy}이지 exact per-patch energy가 아니다.
Mass norm은 \(\|u\|_M:=\sqrt{u^\top Mu}\)로 정의한다. 따라서 displacement numerator의 단위는
\(\sqrt{M}L\), velocity numerator의 단위는 \(\sqrt{M}L/T\)이고, displacement/velocity proxy는
\[
  \widehat d_{x,r}^t
  =
  \frac{\|\Psi_r\widetilde z_r^t\|_M}
       {\sqrt{M_{r,\mathrm{ref}}}L_{\mathrm{ref}}},
  \qquad
  \widehat d_{v,r}^t
  =
  \frac{\|\Psi_r\dot{\widetilde z}_r^{\,t}\|_M}
       {\sqrt{M_{r,\mathrm{ref}}}V_{\mathrm{ref}}}
\]
로 고정한다.
분모에 \(M_{r,\mathrm{ref}}\) 자체를 곱하는 표현은 무차원이 아니고 asset mass scale에 의존하므로
Minimal V0 contract에 포함하지 않는다.
\(\gamma_r^t=0\)이고
\(\widehat E_r^t<\tau_E\), \(\widehat d_{x,r}^t<\tau_x\),
\(\widehat d_{v,r}^t<\tau_v\)와 zeroing 시 예상 anchor pop
\(\delta_{\mathrm{pop},r}^t<\tau_{\mathrm{pop}}\)가
\(N_{\mathrm{hold}}\) frame 연속 성립할 때만 slot을 release하고 zero로 만든다.
Frame \(t\)에서 이 네 proxy와 hold predicate는 식~\eqref{eq:minimal-coupled-step}이 산출한
post-solve, pre-zero \((\widetilde z_r^t,\dot{\widetilde z}_r^t)\)로 평가한다.
여기서 scale-aware pop proxy는
\[
  \delta_{\mathrm{pop},r}^t
  =
  \max_{i\in\operatorname{supp}(\Psi_r)}
  \frac{
    \| (\Psi_r\widetilde z_r^t)_i \|_2
    +\Delta t\,\| (\Psi_r\dot{\widetilde z}_r^{\,t})_i \|_2
  }{L_{\mathrm{ref}}+\varepsilon_L}
\]
로 고정한다.
다시 선택되면 state reset 없이 Active로 돌아간다.
별도 Dormant dynamics, learned hysteresis와 future-benefit label은 만들지 않는다.

\section{Direct generalized projection과 single coupled solve}

\subsection{Selected basis와 reduced blocks}

Admission 직후의 selection/decay union을
\(\mathcal S_t^{\mathrm{solve}}
:=\mathcal A_t^{\mathrm{pre}}\cup\mathcal D_t^{\mathrm{pre}}\)라 한다.
아래 식의 compact subscript \(\mathcal S_t\)는 모두 이 pre-release
\(\mathcal S_t^{\mathrm{solve}}\)를 뜻한다.

\begin{equation}
  B_{\mathcal S_t}
  =
  \left[\Phi,\Psi_{\mathcal S_t}\right].
  \label{eq:minimal-selected-basis}
\end{equation}

\begin{equation}
\begin{aligned}
  M_{\mathcal S_t}&=B_{\mathcal S_t}^\top MB_{\mathcal S_t},\\
  C_{\mathcal S_t}&=B_{\mathcal S_t}^\top CB_{\mathcal S_t},\\
  K_{\mathcal S_t}&=B_{\mathcal S_t}^\top KB_{\mathcal S_t}.
\end{aligned}
  \label{eq:minimal-coupled-blocks}
\end{equation}

모든 cross block을 포함한다. 작은 active rank에서는 precomputed block을 gather해
dense Cholesky/LDL로 직접 푼다. Incremental Cholesky, block-PCG와 dynamic cache는
profiling에서 factorization이 실제 병목으로 확인되기 전에는 구현하지 않는다.

\subsection{Force는 한 번만 투영한다}

Core quantitative run은 gravity-off로 고정하고 \(f_{\mathrm{ext}}^t=f_{\mathrm{aero}}^t\)로 둔다.
따라서 zero-wind rest-recovery fixture의 목표는 undeformed rest다.
Gravity는 Gate A/B/C 뒤 별도 pre-equilibrated sag preset으로만 추가한다.
Global block에는 항상 계수 1을, Active Local에는 \(\gamma_r^t=1\),
Decay Local에는 \(0\leq\gamma_r^t<1\)을 적용하는 block diagonal matrix를
\(\Gamma_{\mathcal S_t}^t\)라 한다. Decay에서 줄이는 것은 새 external excitation뿐이며
structural target, inertia, stiffness와 damping에는 gate를 곱하지 않는다.

\begin{equation}
  r_{\mathcal S_t}^t
  =
  B_{\mathcal S_t}^\top\bar h^t
  +\Gamma_{\mathcal S_t}^t
  B_{\mathcal S_t}^\top f_{\mathrm{ext}}^t.
  \label{eq:minimal-direct-load}
\end{equation}

하나의 external anchor force가 \(\Phi\)와 \(\Psi_{\mathcal S_t}\) coordinate에 동시에 투영되는 것은
double counting이 아니다. 이는 결합 Galerkin basis에서 같은 virtual work를 표현한 것이고,
질량/강성/감쇠 cross block을 포함한 한 system에서 함께 푼다.

Physics-only V0는 full-anchor force를 patch로 잘랐다가 다시 조립하지 않는다.
따라서 nodal lifting, owner-weighted force assembly, exact wrench KKT와 support reaction이 없다.

\subsection{한 번의 implicit midpoint step}

Reduced coordinate를 \(y_{\mathcal S_t}=[q,z_{\mathcal S_t}]\)로 두고,
frame 동안 \(h^t,f_{\mathrm{aero}}^t\)를 hold한다.

\begin{equation}
\begin{aligned}
  A_{\mathcal S_t}^{\mathrm{MP}}
  &=
  \frac{4}{\Delta t^2}M_{\mathcal S_t}
  +\frac{2}{\Delta t}C_{\mathcal S_t}
  +K_{\mathcal S_t},\\
  A_{\mathcal S_t}^{\mathrm{MP}}\Delta y^t
  &=
  2r_{\mathcal S_t}^t
  -2K_{\mathcal S_t}y^{t-1}
  +\frac{4}{\Delta t}M_{\mathcal S_t}\dot y^{t-1},\\
  \widetilde y^t&=y^{t-1}+\Delta y^t,\qquad
  \dot{\widetilde y}^{\,t}
  =\frac{2}{\Delta t}\Delta y^t-\dot y^{t-1}.
\end{aligned}
  \label{eq:minimal-coupled-step}
\end{equation}

Implicit midpoint는 high-frequency visual band를 과도하게 지우지 않으면서
fixed-Hessian solve를 유지하기 위한 구현 선택이며 novelty가 아니다.
Backward Euler는 robustness baseline으로만 비교할 수 있다.
한 frame 안에 aerodynamic force를 다시 평가하거나 Global corrector를 추가하지 않는다.
이 단일 solve 뒤 기존 hold counter를 갱신하고, release predicate를
\(N_{\mathrm{hold}}\) frame 연속 통과한 unit 집합을
\(\mathcal R_t\subseteq\mathcal D_t^{\mathrm{pre}}\)라 한다. Frame-end contract는
\[
  \mathcal D_t^{\mathrm{end}}
  =\mathcal D_t^{\mathrm{pre}}\setminus\mathcal R_t,
  \qquad
  (z_r^t,\dot z_r^t)
  =
  \begin{cases}
    (0,0),&r\in\mathcal R_t,\\
    (\widetilde z_r^t,\dot{\widetilde z}_r^{\,t}),&r\notin\mathcal R_t,
  \end{cases}
  \qquad
  (q^t,\dot q^t)=(\widetilde q^t,\dot{\widetilde q}^{\,t})
\]
이다. Reconstruction은 이 frame-end state를 사용하고, 다음 frame admission은
\(\mathcal D_t^{\mathrm{end}}\)에서 시작한다. Frame \(t\) actual-rank/timing trace는
\(\mathcal S_t^{\mathrm{solve}}\)를 기록한다. 이 순서는 solve를 추가하거나 새 runtime state/stage를 만들지 않는다.

\section{Anchor-to-Gaussian affine transport}

Gaussian \(j\)는 setup에서 고정한 same-layer nearby anchor, weight \(\omega_{ji}\),
rest tangent frame \(T_j^0=[t_{j1}^0,t_{j2}^0]\), normal \(n_j^0\)와
2D rest coordinate \(\xi_{ji}\)를 가진다.
Thin-surface neighborhood만으로 일반 3\(\times\)3 affine map의 normal column을 식별할 수 없으므로,
먼저 tangent map만 fit한다.

Setup binding은 식~\eqref{eq:minimal-surface-tangent-fit}과 동일한
\(\texttt{tangent\_fit\_v1}\)을 사용한다.
즉 positive normalized weights, weighted rest/current centering과
fixed rest coordinate convention을 그대로 따른다.
\begin{equation}
\begin{aligned}
  &\omega_{ji}>0,\qquad
  \sum_{i\in\mathcal N_j}\omega_{ji}=1,\\
  &\bar x_j^s=\sum_{i\in\mathcal N_j}\omega_{ji}x_i^s
  \quad(s\in\{0,t\}),\qquad
  \xi_{ji}=(T_j^0)^\top(x_i^0-\bar x_j^0),\\
  &\sum_i\omega_{ji}\xi_{ji}=0,\qquad
  G_j^0=\sum_i\omega_{ji}\xi_{ji}\xi_{ji}^\top,\qquad
  G_j^0\ \text{passes}\ \eqref{eq:minimal-tangent-numerical-rank},\quad
  \kappa(G_j^0)\leq\kappa_{\mathrm{MLS}}.
\end{aligned}
  \label{eq:minimal-affine-weights}
\end{equation}
따라서 각 support에는 최소 세 개의 non-collinear same-layer anchor가 필요하다.
Rest와 current support에서 같은 \(\xi_{ji}\)에 대한 3\(\times\)2 fitted tangent Jacobian을
각각 구하고, rest fit에 대한 상대 map만 transport에 사용한다.

\begin{equation}
\begin{aligned}
  C_j^s
  &=\sum_{i\in\mathcal N_j}\omega_{ji}
    (x_i^s-\bar x_j^s)\xi_{ji}^\top,\\
  J_j^s
  &=C_j^s(G_j^0)^{-1},\\
  J_j^s
  &=\arg\min_{J\in\mathbb R^{3\times2}}
  \sum_{i\in\mathcal N_j}\omega_{ji}
  \left\|(x_i^s-\bar x_j^s)-J\xi_{ji}\right\|_2^2,
  \qquad s\in\{0,t\},\\
  F_{j,\parallel}^t
  &=J_j^t\bigl((J_j^0)^\top J_j^0\bigr)^{-1}(J_j^0)^\top .
\end{aligned}
  \label{eq:minimal-affine-fit}
\end{equation}

\(J_j^0\)와 \(J_j^t\)는 식~\eqref{eq:minimal-tangent-numerical-rank}을 통과하고
\(\kappa((J_j^s)^\top J_j^s)\leq\kappa_{J,\mathrm{MLS}}\),
\(s\in\{0,t\}\)여야 한다.
Ordered tangent column과 rest orientation sign \(s_j\in\{-1,+1\}\)로
\[
  \widetilde n_j^s=J_j^s(:,1)\times J_j^s(:,2),\qquad
  n_{j,\mathrm{fit}}^s
  =s_j\frac{\widetilde n_j^s}{\|\widetilde n_j^s\|_2},
  \quad s\in\{0,t\},
\]
를 정한다. \(s_j\)는 \(n_{j,\mathrm{fit}}^0\)가 stored \(n_j^0\)와 같은 방향이 되도록
setup에서 동결한다. Rest \(G_j^0,J_j^0\)의 numerical-rank/condition과
orientation 또는 normal norm이 cap을
통과하지 못하면 binding/package를 reject한다. Valid affine path에서는 V0의 declared unit normal
extension을 사용해

\begin{equation}
  F_{j,\mathrm{shell}}^t
  =
  F_{j,\parallel}^t
  +n_{j,\mathrm{fit}}^t(n_{j,\mathrm{fit}}^0)^\top
  \label{eq:minimal-shell-extension}
\end{equation}

을 만든다. 이 상대 정의는 identity에서 \(F_{j,\mathrm{shell}}^0=I\),
proper rigid motion \(Q\)에서 \(F_{j,\mathrm{shell}}^t=Q\)를 정확히 재현한다.

Current \(J_j^t\)의 numerical-rank/condition 또는 normal norm이 cap을 통과하지 못하면
renderer transport consumer는 다음 centered weighted proper-Kabsch fallback을 시도한다.
\begin{equation}
\begin{aligned}
  r_{ji}^s&=x_i^s-\bar x_j^s,\qquad s\in\{0,t\},\\
  H_{j,\mathrm{rigid}}^t
  &=\sum_{i\in\mathcal N_j}\omega_{ji}r_{ji}^t(r_{ji}^0)^\top
    =U_j\operatorname{diag}(\sigma_{j1},\sigma_{j2},\sigma_{j3})V_j^\top,\\
  R_{j,\mathrm{rigid}}^t
  &=U_j\operatorname{diag}
    \left(1,1,\det(U_jV_j^\top)\right)V_j^\top\\
  &=\arg\min_{R\in\mathrm{SO}(3)}
    \sum_{i\in\mathcal N_j}\omega_{ji}
    \left\|r_{ji}^t-Rr_{ji}^0\right\|_2^2 .
\end{aligned}
  \label{eq:minimal-renderer-rigid-fallback}
\end{equation}
Singular values는 \(\sigma_{j1}\geq\sigma_{j2}\geq\sigma_{j3}\geq0\)로 정렬한다.
\(0<\tau_K<1\)에 대해 numerical rank는
\[
  \operatorname{rank}_{\tau_K}(H)
  :=\#\left\{k:\sigma_1>0,\ \sigma_k/\sigma_1\geq\tau_K\right\}
\]
인 relative singular-value convention으로 고정한다.
Fallback 성공은 stored support order/rest orientation에 따른 deterministic SVD sign/tie-break,
\(R_{j,\mathrm{rigid}}^t\in\mathrm{SO}(3)\)와
\[
\begin{aligned}
  \operatorname{rank}_{\tau_K}(H_{j,\mathrm{rigid}}^t)&=2,
  &\frac{\sigma_{j2}}{\sigma_{j1}}&\geq\tau_{K,\mathrm{ratio}},
  &\frac{\sigma_{j3}}{\sigma_{j2}}&\leq\tau_{K,\mathrm{plane}},\\
  \rho_{K,j}^t
  &:=\frac{
    \left(\sum_i\omega_{ji}
    \|r_{ji}^t-R_{j,\mathrm{rigid}}^tr_{ji}^0\|_2^2\right)^{1/2}
  }{
    \left(\sum_i\omega_{ji}\|r_{ji}^0\|_2^2\right)^{1/2}}
  \leq\tau_{K,\mathrm{res}}.
\end{aligned}
\]
를 모두 요구한다. 따라서 reflection, rank-1 collapsed support와 numerical rank-3 ambiguity를
rigid 성공으로 간주하지 않는다. Finite하고 식~\eqref{eq:minimal-tangent-numerical-rank}의 tangent fit과
stored orientation을 통과한 affine candidate의 다음 3\(\times\)3 SVD로 full-rank/orientation 조건을
먼저 검사하고, 통과한 경우에만 singular value를 clip한다.
\begin{equation}
\begin{aligned}
  F_{j,\mathrm{shell}}^t
  &=U_{F,j}^t
    \operatorname{diag}(\sigma_{F,j1}^t,\sigma_{F,j2}^t,\sigma_{F,j3}^t)
    (V_{F,j}^t)^\top,\\
  \sigma_{F,j1}^t&\geq\sigma_{F,j2}^t\geq\sigma_{F,j3}^t
    \geq\sigma_{F,\mathrm{rank}}>0,
  \qquad \det(F_{j,\mathrm{shell}}^t)>0,\\
  \bar\sigma_{F,jk}^t
  &=\operatorname{clip}(\sigma_{F,jk}^t;\sigma_{F,\min},\sigma_{F,\max}),
  \qquad 0<\sigma_{F,\min}\leq1\leq\sigma_{F,\max},\\
  F_{j,\mathrm{safe}}^t
  &=U_{F,j}^t
    \operatorname{diag}(\bar\sigma_{F,j1}^t,
                        \bar\sigma_{F,j2}^t,
                        \bar\sigma_{F,j3}^t)
    (V_{F,j}^t)^\top .
\end{aligned}
  \label{eq:minimal-affine-safe-map}
\end{equation}
여기서 \(\sigma_{F,\mathrm{rank}},\sigma_{F,\min},\sigma_{F,\max}\)는 dimensionless다.
Non-finite value, numerical rank loss 또는 orientation disagreement는
defensive affine-invalid event다. 다음 orientation failure도 같은 event로 처리한다.
\[
  \det(F_{j,\mathrm{shell}}^t)\leq0.
\]
Singular-vector sign을 바꾸어
reflection을 조용히 proper map으로 수선하지 않고, 식~\eqref{eq:minimal-renderer-rigid-fallback}의 기존
proper-rigid fallback을 시도한다. Affine path가 위 식을 통과하면 그 \(F_{j,\mathrm{safe}}^t\)를 사용하고,
affine path는 실패했지만 rigid path가 성공하면
\(F_{j,\mathrm{safe}}^t=R_{j,\mathrm{rigid}}^t\)로 둔다.

Binding package는 \texttt{tangent\_fit\_version}, ordered-neighbor/weight version,
\(\kappa_{\mathrm{MLS}}\), \(\kappa_{J,\mathrm{MLS}}\),
normal floor, rigid-fallback ID와
\(\tau_K,\tau_{K,\mathrm{ratio}},\tau_{K,\mathrm{plane}},\tau_{K,\mathrm{res}}\),
\(\sigma_{F,\mathrm{rank}},\sigma_{F,\min},\sigma_{F,\max}\), deterministic 3D-SVD ID,
covariance eigenvalue bound와
development에서 동결한 renderer fallback-rate cap \(\tau_{\mathrm{rigid}}\)을 저장한다.
Required affine attempt 중 invalid와 singular-value clip event 비율
\(r_{\mathrm{affine,invalid}},r_{\mathrm{affine,clip}}\)을 각각 trace에 기록하고,
event definition, threshold와 package hash를 함께 저장한다.
Run의 required Gaussian--frame transport attempt 수에 대한 successful rigid-fallback event 비율을
\(r_{\mathrm{rigid}}\)로 기록한다. Rigid fallback이 성공하고
\(r_{\mathrm{rigid}}\leq\tau_{\mathrm{rigid}}\)이면 fallback-marked renderer 결과를 사용할 수 있다.
Rigid fit 자체가 실패하거나 cap을 넘으면 식~\eqref{eq:minimal-attempt-success-contract}의
in-envelope method failure다. 이 fallback은 식~\eqref{eq:minimal-surface-tangent-fit}의
surface/aero consumer에는 적용하지 않는다. Gaussian mean/covariance는

\begin{equation}
\begin{aligned}
  \mu_j^t
  &=
  \bar x_j^t+F_{j,\mathrm{safe}}^t(\mu_j^0-\bar x_j^0),\\
  \widetilde\Sigma_j^t
  &=
  F_{j,\mathrm{safe}}^t\Sigma_j^0(F_{j,\mathrm{safe}}^t)^\top
  =Q_j^t\operatorname{diag}(\lambda_{j1}^t,\lambda_{j2}^t,\lambda_{j3}^t)(Q_j^t)^\top,\\
  \bar\lambda_{jk}^t
  &=\operatorname{clip}(\lambda_{jk}^t;\lambda_{\min},\lambda_{\max}),
  \qquad k\in\{1,2,3\},\\
  \Sigma_j^t
  &=
  Q_j^t\operatorname{diag}(\bar\lambda_{j1}^t,\bar\lambda_{j2}^t,\bar\lambda_{j3}^t)(Q_j^t)^\top .
\end{aligned}
  \label{eq:minimal-gaussian-transport}
\end{equation}

로 한 번만 갱신한다. Eigenvalue가 renderer-safe 범위 안이면 clip은 identity이며,
unconditional \(\varepsilon_\Sigma I\)를 더하지 않는다. 범위를 벗어난 경우에만 conditional clamp하고
clamp rate를 보고한다.
Accepted rest Gaussian은
\(\lambda_k(\Sigma_j^0)\in[\lambda_{\min},\lambda_{\max}]\)를 setup에서 먼저 통과해야 하므로
identity deformation에서는 covariance clamp가 작동하지 않는다.
Identity, rigid transform, in-plane affine reproduction과 declared normal extension, SPD가 hard test다.
Planar evidence로 arbitrary 3D affine 또는 thickness deformation을 복원한다고 주장하지 않는다.

Canonical quantitative rendering은
\(\texttt{appearance\_mode=degree0\_v0}\) 하나만 사용한다.
각 Gaussian의 opacity와 degree-0/DC color coefficient는 rest asset 값을 유지하고,
higher-order SH coefficient는 삭제하지 않되 canonical renderer에서 비활성화한다.
따라서 V0 transport가 갱신하는 것은 mean과 covariance뿐이며, 회전하는 surface의 higher-order
view-dependent lobe를 world frame에 고정한 채 appearance 성능으로 주장하지 않는다.
Local rigid frame으로 view direction 또는 SH frame을 analytic하게 회전하는 방식은
Gate A/B/C 이후 별도 appearance tag를 가진 optional ablation이다. 이는 V0 DoD가 아니며
appearance residual network나 새 runtime physics stage를 추가하는 근거가 아니다.

이 transport는 known physics-inspired upsampling과 연결되며
\cite{kavan2011upsampling}, 독립 algorithm novelty로 주장하지 않는다.

\section{Minimal safety contract}

다음 여덟 항목만 구현 차단 hard gate로 둔다.

\begin{longtable}{L{0.08\linewidth}L{0.38\linewidth}L{0.44\linewidth}}
\toprule
ID & Hard acceptance & 실패 시 조치\\
\midrule
\endfirsthead
\toprule
ID & Hard acceptance & 실패 시 조치\\
\midrule
\endhead
S0 & Frame/unit/index, SI/nd typed mapping, operating envelope와 \(N_{\mathrm{sub}}=1\),
canonical method signature 및 \(\texttt{appearance\_mode=degree0\_v0}\)를 포함한
package/state/run hash와 save--load round trip &
Fail-fast. Dataset/run 생성 금지.\\
S1 & Positive area/mass, total ownership, attachment-line eigengap/gauge, versioned relation-to-constraint
mapping, flat-rest planarity/tangential residual cap, valid hard attachment와 close-layer shortcut 없음 &
Package reject, user hint 또는 explicit scaffold fallback.\\
S2 & \texttt{corotated\_tangent\_plane\_v2} one-local-step fixed-Hessian PD package,
\(M\succ0\), projector rest/rigid/normal-response fixture, zero-wind recovery와 bounded 10--30초 rollout &
Preset/backend 수정 또는 package reject. 새 shell feature 자동 추가 금지.\\
S3 & \texttt{generalized\_eigen\_v1} Global basis,
\(\Phi^\top M\Psi_r\approx0\), reachable \(B_{\mathcal S}\) full rank/conditioned, cross block 보존 &
Local basis prune/rebuild. Dynamic remapping 금지.\\
S4 & \texttt{tangent\_fit\_v1}의 fixed same-layer weights, rank/condition/orientation,
3\(\times\)2 tangent normal/area, zero-relative-wind force, finite load와 force channel 1회 소비 &
Frame/package fail. H/corrector로 우회 금지.\\
S5 & Direct \(B_{\mathcal S}^\top f\), one coupled solve, frozen dwell/fade/release,
no energetic reset와 activation/release pop tolerance &
Admission 지연. 반복 실패 시 MC2/selectivity를 철회하고 always-on Local baseline으로 재분류;
canonical selective result로 사용 금지.\\
S6 & Normalized rank-2 affine support, 식~\eqref{eq:minimal-renderer-rigid-fallback}의
renderer-only proper-rigid fallback과 rate cap, identity/rigid/in-plane-affine transport,
declared normal extension, conditional covariance clamp와 SPD &
Affine 실패 후 valid rigid와 cap 이하는 marked render 허용.
Rigid 실패/cap 초과는 method failure; surface/aero fallback 금지.\\
S7 & GPU synchronization/warmup, p50/p95/p99, dynamics/end-to-end 분리,
actual active rank 기록 &
Speed/real-time/actual-skip claim 금지.\\
\bottomrule
\end{longtable}

S0--S7은 engineering fidelity test가 아니라 visible failure와 claim invalidation을 막는 계약이다.
Exact force/torque, exact reaction과 exact modal closure를 hard gate로 다시 올리지 않는다.

\section{Topology distillation contract}

\subsection{V0-R1 versioned executable package}

Topology distillation은 runtime network가 아니라 setup용 V0-R1 module 하나다. 그 실행 계약을
\begin{equation}
  \mathcal V_{\mathrm{R1}}
  =(\nu_{\mathrm{anchor}},\nu_{\mathrm{cand}},\nu_{\mathrm{feat}},
    \nu_{\mathrm{arch}},\nu_{\mathrm{label}},\nu_{\mathrm{loss}},
    \nu_{\mathrm{decode}},\nu_{\mathrm{split}})
  \label{eq:minimal-r1-version-contract}
\end{equation}
로 저장한다. 각 \(\nu\)는 algorithm ID, normalization, threshold/cap, seed와 deterministic tie-break를
포함한다. Surface-valid Gaussian에서 anchor를 만드는
\(\mathcal S_{\mathrm{anchor}}^{\nu_{\mathrm{anchor}}}\)와 각 anchor의 broad proposal set
\(\mathcal P_i^{\nu_{\mathrm{cand}}}\)를 config에서 정확히 하나 선택하고
\begin{equation}
  V=\mathcal S_{\mathrm{anchor}}^{\nu_{\mathrm{anchor}}}(\mathcal G_{\mathrm{valid}}),
  \qquad
  E_{\mathrm{cand}}=\bigl\{\{i,k\}:k\in\mathcal P_i\ \text{or}\ i\in\mathcal P_k\bigr\}
  \label{eq:minimal-r1-anchor-candidate}
\end{equation}
로 undirected candidate graph를 만든다. Weighted FPS, deterministic radius/kNN 또는 이들의 고정
조합은 허용 가능한 config 후보지만 어느 하나를 방법 전체의 유일한 선택으로 선언하지 않는다.
각 training run은 아래 lifecycle에 따라 routine, cardinality/radius와 normalization이 모두 채워진
정확히 한 training-static config만 사용하며, asset별 retuning을 금지한다.

Object center/scale로 정규화한 position, versioned positive Gaussian surface scale
\(\ell_i^{\mathrm{GS}}>0\)와 eigenvalue ratio, surface normal의 sign-invariant
quantity, attachment distance와 식~\eqref{eq:minimal-local-line-gauge}의 line projector로 node feature
\(\phi_i^{\nu_{\mathrm{feat}}}\)를 만든다.
\(\bar\ell_{ik}^{\mathrm{GS}}=(\ell_i^{\mathrm{GS}}+\ell_k^{\mathrm{GS}})/2\),
\(r_{ik}=(x_k^0-x_i^0)/\bar\ell_{ik}^{\mathrm{GS}}\)와
\(\{a,b\}_{\mathrm{sym}}=(a+b,|a-b|)\)를 사용해 pair feature를
\begin{equation}
\begin{aligned}
  \chi_{ik}=\bigl[&\|r_{ik}\|_2,\ |(n_i^0)^\top n_k^0|,
  \{ |r_{ik}^\top n_i^0|,|r_{ik}^\top n_k^0|\}_{\mathrm{sym}},\\
  &\{r_{ik}^\top L_i^0r_{ik},r_{ik}^\top L_k^0r_{ik}\}_{\mathrm{sym}},
  |\log(\ell_i^{\mathrm{GS}}/\ell_k^{\mathrm{GS}})|\bigr]
\end{aligned}
  \label{eq:minimal-r1-symmetric-features}
\end{equation}
로 고정한다. Feature의 정확한 field order와 epsilon은 \(\nu_{\mathrm{feat}}\)가 소유한다.
Node encoder \(h_i=E_\theta(\phi_i)\) 뒤 relation/layer logit은
\begin{equation}
  \ell_{ik}^{c}
  =g_{\theta,c}\!\left(h_i+h_k,|h_i-h_k|,\chi_{ik}\right),\qquad
  p_{ik}^{c}=\sigma(\ell_{ik}^{c})=p_{ki}^{c},\quad
  c\in\{\mathrm{rel},\mathrm{layer}\}
  \label{eq:minimal-r1-symmetric-logit}
\end{equation}
처럼 symmetry를 construction으로 보장한다. Width/depth/activation은
\(\nu_{\mathrm{arch}}\)에서 하나로 동결하되 특정 architecture를 method claim으로 만들지 않는다.
Area head는 node embedding \(h_i\)만 소비하는 scalar head이며 global pooling이나 추가
message-passing stage를 사용하지 않는다. Dimensionless floor
\(\varepsilon_{a,\mathrm{raw}}>0\)에 대해
\begin{equation}
\begin{aligned}
  \zeta_i^a
  &=g_{\theta,a}(h_i),\\
  \widetilde a_i
  &=\operatorname{softplus}(\zeta_i^a)+\varepsilon_{a,\mathrm{raw}},\\
  a_i^0
  &=A_{\mathrm{ref}}\frac{\widetilde a_i}{\sum_{k\in V}\widetilde a_k}.
\end{aligned}
  \label{eq:minimal-r1-area-decode}
\end{equation}
로 \(\lvert V\rvert\times1\) scalar output과 positive total ownership을 construction으로 만족한다.
Area-head width/depth/activation과 output field ID는 \(\nu_{\mathrm{arch}}\)가 소유한다.
여기서 unit-typed \(A_{\mathrm{ref}}>0\)는 물리적 절대 면적 추정값이 아니라
식~\eqref{eq:minimal-area-mass}의 deterministic 소유권 정규화 상수다.
Package build 전에 \(L_{\mathrm{ref}}=\operatorname{diameter}(x^0)>0\),
\(A_{\mathrm{ref}}=L_{\mathrm{ref}}^2\)로 고정하고, object-normalized nondimensional geometry에서
\(L_{\mathrm{ref}}=1\)이면 \(A_{\mathrm{ref}}=1\)이다. Training source mesh나 target static GS에
별도의 absolute-area calibration estimator를 만들지 않는다.

Training-only source mesh에는 anchor-to-surface map과 label generator
\(\mathcal Y^{\nu_{\mathrm{label}}}\)를 적용해 relation, close-layer/shortcut과 area target
\((y_{ik}^{\mathrm{rel}},y_{ik}^{\mathrm{layer}},a_i^\star)\)를 만든다. Connectivity-hop, geodesic-radius
또는 versioned side/layer rule 중 선택한 정확한 한 규칙과 모든 radius/hop/mapping tie-break를 config/hash에
저장하며, geodesic radius 하나를 모든 asset의 유일한 ground truth로 강제하지 않는다. Source mesh는
label 생성에만 쓰고 target setup/runtime 입력에는 들어가지 않는다. 식~\eqref{eq:minimal-attachment-line-gauge}의
in-plane gauge는 analytic이므로 orientation label/output/loss를 만들지 않고 \(\Delta\theta_i=0\)을 유지한다.
Binary label polarity는
\[
\begin{aligned}
  y_{ik}^{\mathrm{rel}}=1
  &\iff \{i,k\}\text{가 valid same-layer material relation이고},\\
  y_{ik}^{\mathrm{layer}}=1
  &\iff \{i,k\}\text{가 invalid cross-layer 또는 geometric shortcut candidate다}
\end{aligned}
\]
로 고정한다. 따라서 valid relation은 \((y_{ik}^{\mathrm{rel}},y_{ik}^{\mathrm{layer}})=(1,0)\)이고,
invalid layer/shortcut은 relation positive로 중복 표기하지 않는다.
\(p_{ik}^{\mathrm{rel}}\)과 \(p_{ik}^{\mathrm{layer}}\)는 각각 이 positive class의 probability이며,
label enum/string, model-output metadata와 decoder metadata를 \(\nu_{\mathrm{label}}\) hash에 함께 저장한다.

Class weight가 동결된 weighted BCE를 \(\operatorname{BCE}_{w}\)라 하면
\begin{equation}
\begin{aligned}
  \mathcal L_{\mathrm{edge}}
  &=|E_{\mathrm{cand}}|^{-1}\sum_{\{i,k\}}\operatorname{BCE}_{w_e}
    (p_{ik}^{\mathrm{rel}},y_{ik}^{\mathrm{rel}}),\\
  \mathcal L_{\mathrm{layer}}
  &=|E_{\mathrm{cand}}|^{-1}\sum_{\{i,k\}}\operatorname{BCE}_{w_l}
    (p_{ik}^{\mathrm{layer}},y_{ik}^{\mathrm{layer}}),\\
  \mathcal L_{\mathrm{area}}
  &=|V|^{-1}\sum_i\left[\log(a_i^0/A_{\mathrm{ref}}+\varepsilon_{a,\log})-
    \log(a_i^\star/A_{\mathrm{ref}}+\varepsilon_{a,\log})\right]^2,\\
  \mathcal L_{\mathrm{topo}}
  &=\lambda_e\mathcal L_{\mathrm{edge}}
    +\lambda_a\mathcal L_{\mathrm{area}}
    +\lambda_{\mathrm{layer}}\mathcal L_{\mathrm{layer}}.
\end{aligned}
  \label{eq:minimal-topology-loss}
\end{equation}
여기서 dimensionless \(\varepsilon_{a,\log}>0\)를 사용하고 teacher area target도
합이 \(A_{\mathrm{ref}}\)가 되도록 정규화한다.
Loss weight, 두 dimensionless epsilon, optimizer/seed와 stopping rule은 \(\nu_{\mathrm{loss}}\)에 저장한다.
Frozen decoder는 \(p_{ik}^{\mathrm{rel}}\geq\tau_{\mathrm{rel}}\),
\(p_{ik}^{\mathrm{layer}}\leq\tau_{\mathrm{layer}}\), degree/radius/layer hard check와 stable-ID tie-break를
모두 통과한 undirected edge만 accept한다. Degree cap 충돌은 score와 stable ID로 결정하고 hidden graph
repair나 per-asset threshold 변경을 금지한다.

Split의 독립 단위는 source object group이다. 같은 source object의 GS density/seed/reconstruction variant는
한 group에만 속하고 train/development/test 사이에 나뉘지 않는다. V0-R1 lifecycle은 다음 순서만 허용한다.
\begin{enumerate}
  \item 먼저 \(\nu_{\mathrm{split}}\), object-group assignment와 data hash를 동결하고 test object/package ID를
  seal한다. 이후 config/model/decoder 선택 동안 test payload와 result를 열지 않는다.
  \item 각 training entry 전에
  \(\nu_{\mathrm{anchor}},\nu_{\mathrm{cand}},\nu_{\mathrm{feat}},\nu_{\mathrm{arch}},
  \nu_{\mathrm{label}},\nu_{\mathrm{loss}},\nu_{\mathrm{split}}\), decoder schema와 threshold/cap
  candidate rule, seed 및 training data/label hash가 non-placeholder인 training-static config인지 검사하고
  하나라도 비었으면 fail-fast한다. 이때 최종 numerical threshold/cap은 아직 동결하지 않는다.
  \item Preregistered candidate run은 train group만 학습하고 development group만으로 정확히 한
  model/config와 decoder threshold/cap을 선택한다. Asset별 또는 test-aware retuning은 금지한다.
  \item Test payload를 처음 열기 전에 \(\mathcal V_{\mathrm{R1}}\)의 모든 field와 선택한 model weight를
  동결하며, development에서 선택한 최종 threshold/cap을 포함하는 \(\nu_{\mathrm{decode}}\)도 이 시점에
  처음 완결한다. Training source, label, data, config, decoder와 split의 full hash도 같은 immutable
  manifest에 기록한다.
  Test failure 뒤 어느 field라도 바꾸면 새 preregistered run이지 같은 result의 보정이 아니다.
\end{enumerate}
따라서 training 전 freeze는 decoder schema와 candidate rule의 재현성을 소유하고,
최종 \(\nu_{\mathrm{decode}}\) numerical threshold/cap은 development-only 선택 뒤 test-open 전에 동결한다.
최종 \(\nu_{\mathrm{decode}}\)를 training 전에 요구하거나 test 결과를 보고 다시 고르는 두 순서 모두 금지한다.
이 lifecycle의 미완료나 실패는 V0-R1 predicted scaffold와 Gate A/MC1 evidence만 막는다.
Explicit Oracle scaffold의 V0-01--04, privileged-teacher-ranked Local을 쓰는 Gate B와
dynamics-only Gate C 경로를 막지 않는다.
\(\mathcal V_{\mathrm{R1}}\), training source/label hash, model/config hash, decoder와 accepted/rejected reason을
package에 저장한다. Network가 material, attachment,
dense \(M,C,K\), dynamic state, force 또는 constraint triplet을 예측하지 않는다. Accepted output은
식~\eqref{eq:minimal-topology-to-constraints}의 단일 analytic decoder만 거치며 full confidence calibrator를
추가하지 않는다. Edge AUPRC는 diagnostic이고, 최종 MC1 evidence는 predicted scaffold가 Oracle scaffold에 근접하고
oriented-kNN보다 response robustness와 \(R_{\mathrm{accept}}\)에서 낫다는 것이다.

\subsection{Gate A가 실패할 때}

Confidence network를 추가하는 대신 다음 순서로 축소한다.

\begin{enumerate}
  \item 지원 asset class를 제한한다.
  \item close-layer/user orientation hint를 받는다.
  \item explicit target scaffold 또는 Oracle scaffold fallback을 공개한다.
  \item 그래도 실패하면 MC1을 철회하고 MC2를 Oracle scaffold에서만 진단한다.
\end{enumerate}

\section{연구 Gate A/B/C}

\subsection{Gate A --- Representation}

Independent GS 2--3개씩의 strip/flag에서 Oracle/source scaffold와 predicted scaffold를 비교한다.
필수 증거는 area/mass consistency, close-layer false edge, coarse low-frequency band,
동일 wind의 tip/velocity/PSD response consistency와 density/seed robustness다.
첫 8--16 mode의 exact match는 diagnostic이지 universal hard threshold가 아니다.
식~\eqref{eq:minimal-attempt-success-contract}의 \(R_{\mathrm{accept}}\)도 primary evidence이며,
동결한 \(\tau_{\mathrm{fail},A}\)를 넘는 package failure는 accepted package의 conditional
response가 좋아도 Gate A 실패다.

Gate A 실패는 MC1을 축소하지만 Oracle scaffold에서 Gate B를 자동 기각하지 않는다.

\subsection{Gate B --- Local necessity}

동일 p95 dynamics latency에서

\[
\begin{aligned}
  &\text{strong Global rank sweep}\qquad\text{vs.}\\
  &\text{Global + privileged-teacher-ranked}\\
  &\text{complementary Local}
\end{aligned}
\]

을 비교한다. Privileged-teacher-ranked Local이 tip flutter, traveling gust, velocity trajectory 또는
target-band PSD에서 식~\eqref{eq:minimal-pareto-effects}의 clustered effect-margin improvement를
반복적으로 만들지 못하면
MC2와 selector 개발을 중단하고 Global-only 방향으로 축소한다.

Privileged-teacher-ranked selector는 setup에서 동결한 동일 atomic units와 \(K_A,K_D\),
동일 coupled solver를 사용하되,
privileged teacher의 다음 \(H\) frame velocity/target-band error 감소량으로 unit을 rank한다.
Future label search 비용은 runtime latency에서 제외하지만 선택 뒤 실제 execution 비용은 포함한다.
\(H\), error metric과 tie-break를 development split에서 한 번 동결한다.
Offline scoring protocol은 \texttt{teacher\_counterfactual\_v1}로 고정한다. 각 scoring 시점의 full serialized
\texttt{RuntimeState}를 candidate별로 deep-copy하되 Global/Local coordinate와 velocity,
Active/Decay membership, dwell age, fade/release 상태 및 모든 timer/counter를 빠짐없이 복제한다.
Gate B manifest는 이 protocol ID, \(H\), error metric, tie-break, scoring 시작 initial serialized
\texttt{RuntimeState} hash, wind hash, reference hash와 solver/precision/seed를 포함한 config hash를 소유한다.
Scoring 시작 state에서 minimum dwell을 아직 채우지 않아 반드시 유지해야 하는 Active unit 집합을
\(\mathcal A_{\mathrm{mand}}\)라 한다. 각 candidate \(r\)는 동일 initial copy에서 base desired request
\(\mathcal A_{\mathrm{mand}}\)와 isolated-candidate desired request
\(\mathcal A_{\mathrm{mand}}\cup\{r\}\)를 비교하며, \(r\)은 optional priority 중 최상위로 둔다.
두 branch는 동일 wind/reference horizon, solver, precision과 seed를 사용하고, 각각의 desired-set request를
\(H\) frame 동안 고정한 채 기존 capacity-\(K_A\) admission/fade/release policy로 membership을 resolve한다.
따라서 \(|\mathcal A_{\mathrm{mand}}|=K_A\)이면 candidate가 즉시 admit되지 않을 수 있으며 그 이득은
첫 frame shortcut이 아니라 이후 dwell/admission dynamics를 포함한 전체 horizon으로 판정한다.
중간 state reset, recursive re-ranking과 subset search는 금지한다. 이 counterfactual은 ranking-only이며,
선택 뒤 실제 teacher-ranked execution만 기존 fixed \(K_A\)와 \(K_D\)를 따른다. Candidate가 첫 frame에
admit되지 않았다는 이유만으로 score를 0으로 두지 않고, 두 \(H\)-frame rollout의 전체 state/metric trace가
동일할 때만 0으로 둔다.
이는 per-unit ranking으로 Local 필요성을 빠르게 판정하는 scope-control diagnostic이며,
cross-unit coupling과 redundancy를 탐색하는 exact best-subset oracle이나 formal upper bound가 아니다.

\subsection{Gate C --- Actual deployable gain}

Deterministic selector를 포함한 실제 V0가 다음 중 하나를 보여야 한다.

\begin{itemize}
  \item optimized same-anchor sparse solver 대비 dynamics-only speedup
  \item 동일 p95에서 더 낮은 trajectory/velocity/spectrum error
  \item matched quality의 multi-object batch에서 same-anchor baseline보다 높은 throughput
\end{itemize}

Rendering Gaussian 수 증가와 거의 분리된 dynamics cost는 SYS scaling evidence이며,
그 결과만으로 Gate C 또는 MC2를 통과시키지 않는다.
Final deployable 판정은 식~\eqref{eq:minimal-attempt-success-contract}의
\(R_{\mathrm{success}}\)와 \(\tau_{\mathrm{fail},C}\)도 동시에 통과해야 한다.
In-envelope failure가 cap을 넘으면 성공 sequence의 Pareto가 좋아도 Gate C를 통과하지 않는다.

임의의 \(2\times\)는 보편 합격선이 아니다.
그러나 speed를 headline으로 쓰려면 strong-baseline Pareto envelope에
식~\eqref{eq:minimal-pareto-effects}의 practical margin과 uncertainty를 통과한 새 영역이 있어야 한다.

\section{최소 실험과 baseline}

\subsection{첫 sweep}

\begin{itemize}
  \item Strip과 flag, asset별 independent GS 2--3개
  \item Material preset 2--3개, hard edge attachment 하나
  \item Steady, start/stop, sinusoidal, traveling gust
  \item Curated subset만 사용하고 full Cartesian product를 만들지 않음
\end{itemize}

\subsection{필수 baseline}

\begin{enumerate}
  \item High-resolution prescribed-wind structural reference: accuracy reference 전용
  \item 동일 constraint와 timestep을 쓰는 optimized full-coordinate same-anchor PD backend
  \item Global ROM rank sweep
  \item Global + always-on Local upper bound
  \item Global + privileged-teacher-ranked Local
  \item Global + deterministic selected Local
  \item Oracle scaffold, predicted scaffold, simple oriented-kNN scaffold
  \item 최종 단계의 explicit target-mesh extraction 한 방법
  \item Paper evidence 단계의 direct 3DGS wind/animation baseline 한 방법
\end{enumerate}

초기 kill test에서 여러 sparse solver, 다수 neural dynamics baseline,
MPM speed comparison, H/F/G variant와 corrector sweep을 구현하지 않는다.
Direct 3DGS baseline은 Gate A/B/C의 초기 판정을 막지 않으며,
Gate 뒤 공개 코드와 비교 가능성이 가장 높은 한 방법만 선택해 adaptation contract를 동결한다.
\emph{Gaussian Swaying}은 surface-aware Gaussian aerodynamics가 직접 겹치는 paper-stage 우선
후보지만 확정 baseline은 아니다. 공개 구현, 입력/출력 적합성과 공정한 adaptation 비용을 확인한 뒤
V0-05에서 한 방법만 확정하며 후보 이름을 모두 구현할 의무는 없다.

\subsection{Reference, sequence metric과 clustered inference}

High-resolution reference는 단순히 더 많은 anchor를 쓴 run이 아니다. 모든 method/refinement에
동일하게 적용하는 common evaluation-anchor lumped mass를 \(M_{\mathrm{eval}}\)로 고정한다. 두 refinement
\((h,\Delta t)\), \((h/2,\Delta t/2)\)를 common evaluation anchor와 time grid에 resample하고
\begin{equation}
  e_{\mathrm{ref}}
  =
  \left[
  \frac{\sum_t\|v_{\mathrm{ref}}^{h,\Delta t}(t)
    -v_{\mathrm{ref}}^{h/2,\Delta t/2}(t)\|_{M_{\mathrm{eval}}}^2}
       {\sum_t\|v_{\mathrm{ref}}^{h/2,\Delta t/2}(t)\|_{M_{\mathrm{eval}}}^2+\varepsilon_v}
  \right]^{1/2}
  \leq\tau_{\mathrm{ref}}
  \label{eq:minimal-reference-convergence}
\end{equation}
를 통과해야 accuracy reference로 사용한다. 실패는 method failure로 바꾸지 않고
reference-construction failure와 denominator를 별도 보고하며, 해당 accuracy comparison을 성공 사례처럼
선택적으로 남기지 않는다.

Frame은 독립 표본이 아니라 sequence metric의 내부 sample이다. Sequence \(s\)에 대해
\begin{equation}
\begin{aligned}
  e_{v,s}
  &=
  \left[
  \frac{\sum_t\|v_s^t-v_{\mathrm{ref},s}^t\|_{M_{\mathrm{eval}}}^2}
       {\sum_t\|v_{\mathrm{ref},s}^t\|_{M_{\mathrm{eval}}}^2+\varepsilon_v}
  \right]^{1/2},\\
  e_{\mathrm{tip},s}
  &=
  \left[
  \frac{\sum_t\|x_{\mathrm{tip},s}^t-x_{\mathrm{tip,ref},s}^t\|_2^2}
       {\sum_t\|x_{\mathrm{tip,ref},s}^t-x_{\mathrm{tip,ref},s}^0\|_2^2
        +\varepsilon_x}
  \right]^{1/2}
\end{aligned}
  \label{eq:minimal-sequence-errors}
\end{equation}
를 계산한다. 이 절의 \(f_{\mathrm{PSD}}\)는 temporal frequency다.
여기서 canonical lumped \(M_{\mathrm{eval}}\)이 mass-typed일 때
\[
\begin{aligned}
  &\varepsilon_v>0\ \text{and finite},
  &&[\varepsilon_v]=ML^2/T^2,
  &&\widehat\varepsilon_v
    =\varepsilon_v\frac{T_0^2}{M_0L_0^2},\\
  &\varepsilon_x>0\ \text{and finite},
  &&[\varepsilon_x]=L^2,
  &&\widehat\varepsilon_x=\frac{\varepsilon_x}{L_0^2}.
\end{aligned}
\]
역변환은
\(\varepsilon_v=\widehat\varepsilon_vM_0L_0^2/T_0^2\),
\(\varepsilon_x=\widehat\varepsilon_xL_0^2\)로 고정한다.
두 typed value, SI/nondimensional branch, conversion과 \(M_{\mathrm{eval}}\) convention을
evaluation hash에 포함하며, 식~\eqref{eq:minimal-reference-convergence}와
식~\eqref{eq:minimal-sequence-errors}는 같은 \(\varepsilon_v\) contract를 공유한다.
PSD는 per-bin power가 아니라 one-sided velocity-density PSD convention으로 고정하며
\([P]=L^2/T\)다. Sampling rate, window normalization, DC/Nyquist 처리와 development에서 동결한
frequency bin \(B\)를 evaluation hash에 저장한다. PSD와 같은 unit을 가진
\(\varepsilon_P>0\)를 사용하고, unit-parity를 나타내는 기존 typed conversion scale
\((L_0,T_0)\)에 대해
\begin{equation}
\begin{aligned}
  \widehat f_{\mathrm{PSD}}&=f_{\mathrm{PSD}}T_0,&
  \widehat\varepsilon_P&=\varepsilon_P\frac{T_0}{L_0^2},\\
  \widehat P_s(\widehat f_{\mathrm{PSD}})
  &=P_s(f_{\mathrm{PSD}})\frac{T_0}{L_0^2},\\
  \widehat P_{\mathrm{ref},s}(\widehat f_{\mathrm{PSD}})
  &=P_{\mathrm{ref},s}(f_{\mathrm{PSD}})\frac{T_0}{L_0^2},\\
  \delta_s(\widehat f_{\mathrm{PSD}})
  &=\log\!\left(\widehat P_s(\widehat f_{\mathrm{PSD}})+\widehat\varepsilon_P\right)
    -\log\!\left(\widehat P_{\mathrm{ref},s}(\widehat f_{\mathrm{PSD}})+\widehat\varepsilon_P\right),\\
  e_{\mathrm{PSD,RMS},s}
  &=\left(\frac{1}{|\widehat B|}
    \sum_{\widehat f_{\mathrm{PSD}}\in\widehat B}
    \delta_s(\widehat f_{\mathrm{PSD}})^2\right)^{1/2},\\
  e_{\mathrm{PSD,L1},s}
  &=\frac{1}{|\widehat B|}
    \sum_{\widehat f_{\mathrm{PSD}}\in\widehat B}
    |\delta_s(\widehat f_{\mathrm{PSD}})|.
\end{aligned}
  \label{eq:minimal-log-psd-errors}
\end{equation}
를 둘 다 보고한다. Gate와 표의 primary PSD metric은 development에서 두 값 중 하나를 사전 등록하고,
다른 하나는 mandatory secondary로 유지하며 결과를 본 뒤 primary를 바꾸지 않는다.
여기서 \(\widehat B=\{f_{\mathrm{PSD}}T_0:f_{\mathrm{PSD}}\in B\}\)다. 같은 값은
\(\delta_s(f_{\mathrm{PSD}})=
\log[(P_s(f_{\mathrm{PSD}})+\varepsilon_P)/
(P_{\mathrm{ref},s}(f_{\mathrm{PSD}})+\varepsilon_P)]\)로 직접
계산할 수 있으므로 SI direct branch에 새 독립 \(T_{\mathrm{PSD,ref}}\)를 만들지 않는다.
\(\varepsilon_P\)의 typed value/conversion과 위 convention 전체를 evaluation hash에 포함한다.
Signed mean log-PSD처럼 cancellation이 가능한 값 하나로 대체하지 않는다.

같은 source object의 모든 reconstruction은 하나의 outer object cluster에 속하고,
각 independent-GS reconstruction은 그 안의 nested package cluster다. Method--baseline paired difference는
sequence에서 계산하되 uncertainty는 object를 먼저 resample하고, 그 안의 package와 sequence를
순서대로 resample하는 clustered paired bootstrap으로 구한다. Object/package 수가 작은 경우 p-value를
headline으로 쓰지 않고 모든 object/package/sequence effect와 descriptive interval을 공개한다.
Aggregate, bootstrap seed/replicate, failure handling과 cluster 내 weighting은 결과를 보기 전에 동결한다.
독립 source-object 수의 최소값 \(N_{\mathrm{obj,min}}\)도 development에서 사전 등록한다.
그보다 적으면 interval과 Gate 판정을 exploratory/descriptive로 표시하고 accept-ready 통계 주장을 하지 않는다.
표본 부족을 본 뒤 asset을 소급 추가하지 않으며, 다음 paper-evidence 단계의 별도 수집 계획으로만 넘긴다.

Pre-registered baseline-envelope matching rule로 짝지은 method \(m\), baseline \(b\)에 대해
\begin{equation}
  g_{E,s}=\frac{E_{b,s}-E_{m,s}}{E_{b,s}+\varepsilon_{\mathrm{err}}},\qquad
  g_{L,s}=\frac{L_{b,s}-L_{m,s}}{L_{b,s}+\varepsilon_{\mathrm{lat}}}
  \label{eq:minimal-pareto-effects}
\end{equation}
를 사용한다. \(E\)는 위 dimensionless sequence error 중 사전 지정한 primary이고
\(\varepsilon_{\mathrm{err}}>0\)도 dimensionless다. \(L\)은 synchronized sequence latency aggregate이고
\(\varepsilon_{\mathrm{lat}}>0\)는 같은 time unit을 가진다.
Development에서 practical margin \((\delta_E,\delta_L)\)와 non-inferiority tolerance
\((\delta_E^{\mathrm{NI}},\delta_L^{\mathrm{NI}})\)를 동결한다. Pareto improvement는 clustered 95\% interval의
lower bound가
\[
  \operatorname{LB}(g_E)\geq\delta_E,\quad
  \operatorname{LB}(g_L)\geq-\delta_L^{\mathrm{NI}}
  \quad\text{또는}\quad
  \operatorname{LB}(g_L)\geq\delta_L,\quad
  \operatorname{LB}(g_E)\geq-\delta_E^{\mathrm{NI}}
\]
를 만족할 때만 primary evidence다. Margin 없는 strict non-dominance는 diagnostic으로만 남긴다.

\subsection{지표}

\begin{longtable}{L{0.22\linewidth}L{0.30\linewidth}L{0.38\linewidth}}
\toprule
목적 & Primary & 필수 보조\\
\midrule
\endfirsthead
\toprule
목적 & Primary & 필수 보조\\
\midrule
\endhead
Gate A & GS variant response consistency와 \(R_{\mathrm{accept}}\) &
area/mass, false edge, coarse spectrum, failure reason\\
Gate B & velocity error 대 p95 dynamics & tip trajectory, dominant frequency/PSD\\
Gate C & p95 dynamics/throughput와 \(R_{\mathrm{success}}\) &
p99, memory, end-to-end latency, time-to-failure\\
Safety & zero-wind recovery와 in-envelope method-failure rate &
activation jerk/pop, invalid covariance/fallback rate\\
SYS & dynamics cost 대 Gaussian 수 & transport/render time 분리\\
\bottomrule
\end{longtable}

Teacher error가 낮더라도 same-anchor wall-clock가 나쁘면 speed contribution으로 보지 않는다.
반대로 exact force error가 커도 silhouette, trajectory, dominant band와 temporal coherence가
설득력 있고 Pareto가 좋으면 CG 목적에는 허용한다.

\section{Minimal V0 method/contract freeze와 deferred-module policy}

위 canonical signature와 V0-R1 version contract가 이 method/contract freeze를 닫는다.
Gate A/B/C evidence 전에는 명백한 수식 오류의 정정 외에 새 method equation, backend 또는 runtime stage를
추가하지 않는다. 정정은 기존 signature의 의미를 바꾸지 않을 때만 허용하며 변경 hash를 기록한다.

\begin{longtable}{L{0.22\linewidth}L{0.67\linewidth}}
\toprule
Deferred module & 재검토 조건\\
\midrule
\endfirsthead
\toprule
Deferred module & 재검토 조건\\
\midrule
\endhead
H: aero hyper-reduction & Full-anchor aero가 실제 synchronized p95 dynamics의 지배적 병목일 때\\
F: missing-force network & Physics-only Local에 반복 가능하고 시각적으로 중요한 defect가 남을 때\\
G: learned selector & Teacher-ranked와 deterministic selector의 same-\(K\)/same-latency gap이 클 때\\
KKT/Schur & F가 서로 겹치는 nodal patch force를 실제 출력하도록 채택될 때\\
Same-frame corrector & One-frame lag가 visible phase jump/instability를 만들고 same-p95에서 개선될 때\\
Confidence calibrator & Hard triage의 reject/fallback rate가 MC1을 실제로 막을 때\\
Teacher-FSI & Wake/unsteady-aero claim을 논문에 추가하기로 결정할 때\\
Advanced solver backend & Direct small dense solve의 factorization이 profiler에서 병목일 때\\
\bottomrule
\end{longtable}

새 module은 core benchmark에서 다음 중 하나를 실제로 보여야만 승격한다.

\begin{enumerate}
  \item 발산 또는 눈에 띄는 visual artifact 제거
  \item MC1/MC2를 직접 무효화하는 실패 해결
  \item 동일 p95에서 유의미한 quality 개선
  \item matched quality에서 유의미한 latency/throughput 개선
\end{enumerate}

이론적 edge case만 있고 core 결과에 영향이 없으면 limitation으로 종료한다.
새 core stage를 추가하면 기존 stage 하나를 대체해야 하며 병렬 mainline을 만들지 않는다.

\section{Claim 축소와 실패 경로}

\begin{longtable}{L{0.17\linewidth}L{0.17\linewidth}L{0.57\linewidth}}
\toprule
MC1 & MC2 & 연구 경로\\
\midrule
\endfirsthead
\toprule
MC1 & MC2 & 연구 경로\\
\midrule
\endhead
Pass & Pass & Topology-distilled scaffold + selective complementary wind dynamics\\
Pass & Fail & Topology-distilled reduced physical package/Global animation 중심으로 축소\\
Fail & Pass & Oracle/explicit scaffold 기반 selective ROM; novelty/venue 재검토\\
Fail & Fail & 현재 주제 중단 또는 mesh-proxy/visual wind controller로 재설계\\
\bottomrule
\end{longtable}

Selector의 error relevance가 증명되기 전 제목은 \emph{Selective}로 유지한다.
Compliance score가 유효하면 \emph{Compliance-Guided}를 사용할 수 있지만,
certified error estimator 또는 Error-Triggered라고 주장하지 않는다.
현재 canonical 제목의 \emph{Wind Dynamics}는 작업용 framing이다. Prescribed one-way wind에 의해
구동된다는 경계를 더 정확히 드러내는 \emph{Wind-Driven Dynamics} 변경은 Gate A/B/C와
식~\eqref{eq:minimal-pareto-effects}의 effect evidence를 확인한 V0-05에서만 결정한다.
Gate 결과 전에는 특정 venue의 accept 가능성을 claim으로 쓰지 않고 venue 선택도 V0-05 이후로 보류한다.

\section{구현 순서와 Definition of Done}

\subsection{구현 순서}

\begin{enumerate}
  \item \textbf{V0-00:} typed contract, complete SI/nondimensional map,
  operating envelope, attempt/failure 분모, method/contract freeze와 S0 fixture
  \item \textbf{V0-01:} attachment-line gauge, Oracle flat-rest strip scaffold,
  fixed edge/triplet PD backend와 corotated tangent-plane projector,
  topology-to-constraint mapping, \texttt{generalized\_eigen\_v1} Global basis,
  patch-local complement와 S1--S3
  \item \textbf{V0-02:} full-anchor aero, Global rank sweep, same-anchor baseline과 S4
  \item \textbf{V0-03:} privileged-teacher-ranked Local, deterministic selector,
  frozen decay/release schedule,
  single coupled solve와 S5
  \item \textbf{V0-04:} shared tangent-fit contract,\newline
  normalized/rank-checked affine transport,
  renderer-only proper-rigid fallback와 cap, degree-0 appearance,
  conditional covariance clamp,
  integrated profiling과 S6--S7, transport/render를 포함한 final deployable Gate C
  \item \textbf{V0-R1:} strip/flag independent-GS topology distillation과 Gate A
  \item \textbf{V0-05:} Gate A/B/C 결과 집계 및 claim/research-route freeze decision
\end{enumerate}

V0-R1은 V0-01에서 scaffold 요구가 고정된 뒤 V0-02--04와 병렬 수행할 수 있다.
각 P0 contract는 위에서 소유한 milestone 전에만 닫으면 되며,
뒤 milestone의 transport/state 세부사항 때문에 앞선 Oracle fixture 구현을 막지 않는다.
Gate B는 V0-03의 privileged-teacher-ranked Local 결과가 생기면,
Gate C의 dynamics-only 판정은 V0-02 baseline과
V0-03 deterministic timing이 생기면 즉시 먼저 실행할 수 있다.
Transport/render를 포함한 final deployable Gate C는 V0-04 직후 실행한다.
V0 method/contract와 architecture는 구현 전에 이미 위 canonical source로 동결한다.
V0-05는 early evidence를 종합해 최종 claim과 research route를 동결하는 단계이지
판정을 강제로 늦추는 barrier가 아니다.
기존 full-design implementation/experiment 완료 상태는 자동 승계하지 않는다.

\subsection{Definition of Done}

\begin{itemize}
  \item Oracle 및 accepted predicted package에서 관련 S0--S7 통과
  \item Gate A/B/C 각각 pass/fail/defer와 claim 축소 경로 기록
  \item Runtime trace에 deferred module이 존재하지 않음
  \item Strip/flag 핵심 영상과 failure case를 공개할 수 있음
  \item \nolinkurl{basis_type}, \nolinkurl{tangent_fit_version},
  \nolinkurl{tangent_gauge_version}, \nolinkurl{bend_projection_version},
  \nolinkurl{core_scope}와
  derived array/report를 포함한 package hash가 run manifest와 일치
  \item Degree-0 appearance mode, metric/Rayleigh mapping과 renderer fallback
  threshold/cap을 run hash에서 재현 가능
  \item Command, hardware, precision, unit branch와 operating envelope 재현 가능
  \item Pre-registered denominator, \(R_{\mathrm{accept}}\), \(R_{\mathrm{success}}\),
  first-failure reason과 conditional quality/timing을 함께 보고
  \item Dynamics, transport와 render p50/p95/p99를 분리 보고
\end{itemize}

\section{최종 Minimal V0 요약}

Minimal V0는 static GS에서 topology-distilled sparse scaffold를 만들고,
flat-rest core에서 one-local-step fixed-Hessian PD-style structural backend,
\texttt{generalized\_eigen\_v1} Global basis와 full-anchor quasi-steady aero를 사용한다.
Cheap deterministic selector가 선택한 complementary Local basis와 Global basis를
한 coupled Galerkin system으로 한 번 적분한 뒤, aero surface와 Gaussian transport가 공유하는
\texttt{tangent\_fit\_v1} convention으로 Gaussian을 affine transport하고
\texttt{degree0\_v0} appearance로 정량 렌더한다.

방법의 성공은 수식의 완전성이 아니라 다음 세 증거로 판정한다.

\begin{enumerate}
  \item independent GS에서 stable and controllable scaffold를 반복 생성하는가?
  \item strong Global rank보다 selected Local이 같은 p95에서 유용한 detail을 복원하는가?
  \item optimized same-anchor solver를 포함한 Pareto envelope에 새 영역을 만드는가?
\end{enumerate}

이 세 질문에 답하기 전에는 KKT/Schur, corrector, learned residual/gate,
confidence calibration과 새로운 shell mechanics를 추가하지 않는다.

\bibliographystyle{plain}
\bibliography{refs_topology_distilled_selective_global_local_wind_dynamics}

\end{document}
